% ============================================================================
% Chapter 4 — Methodology and Design
% Status: SKELETON. Rubric-mandated two-section top level (4.1 and 4.2);
% depth lives in subsections of 4.2.
%
% Course outcome covered: CO5 (design alternatives, theoretical proofs).
% Preserves the high-level system architecture figure from the earlier
% draft as figures/system_architecture_figure.tex.
%
% Proof discipline: theorems are stated and proved here, next to their
% conceptual justification. Any single proof that grows past half a page
% during writing must be replaced with a proof sketch in the body and
% the full proof relocated to Appendix G. This keeps the chapter
% readable while preserving the theorems as visible contributions.
% ============================================================================

\chapter{Methodology and Design}
\label{ch:methodology}

% Chapter scope: presents the design process and the full architectural
% specification of the system. Section 4.1 walks the reader through the
% methodology at a high level. Section 4.2 unpacks every component of
% the architecture, the algorithms it executes, and the load-bearing
% theorems and corollaries with their proofs, plus the scaffolding
% identities collected as remarks. Each design subsection follows the
% convention: concept, formal definition, rationale for this design
% over alternatives, parameters and how they are selected, complexity,
% and how the choice contributes to the hypotheses registered in
% Chapter 1.

\section{Design Process or Methodology Overview}
\label{sec:design-process}

\subsection{The eight-stage pipeline narrative}
\label{sec:overview-pipeline}

The architecture organises a single query into an eight-stage pipeline that is traversed once per query and wrapped in a per-cycle update loop. Stage~1, pre-routing, computes a calibrated confidence on the query from a single short forward pass that fuses a short-greedy-decode token-probability aggregate with an encoder-layer convergence ratio; the design and rationale appear in \Cref{sec:pre-routing}. Stage~2, encode and lookup, encodes the query into a sentence-embedding vector and probes the episodic memory for the nearest stored entry, returning a similarity score and the neighbour's stored-quality score; the substrate is detailed in \Cref{sec:gen-retrieval} and \Cref{sec:memory-store}. Stage~3, route, combines the similarity and stored-quality scores into a single dispatch score and selects one of three tiers, with an independent safety override that forces dispatch off the score-formula path whenever the pre-routing confidence falls below a fixed floor; the router is specified in \Cref{sec:router}. Stage~3b, retrieval-utility classifier consultation, runs only when the dispatch lands on the retrieval-augmented tier and refines that default into either the zero-shot tier or the retrieval-augmented tier under a learned binary classifier on the query; the classifier is specified in \Cref{sec:ruc}. Stage~4 executes the tier chosen by Stage~3 or refined by Stage~3b: the memory-direct tier serves the neighbour's stored answer; the zero-shot tier runs a generation on the fine-tuned generator; the retrieval-augmented tier retrieves passages from the public corpus and runs a grounded generation. Stage~5, verify, passes every generated answer through the nine-signal post-generation verifier whose four signal families feed the calibrated probability composite; the verifier is specified in \Cref{sec:verifier} and the composite in \Cref{sec:composite}. Stage~6, decide, applies the four-outcome decision tree to the composite output, with a confabulation early-exit rule that takes precedence over the composite thresholds; the decision tree is specified in \Cref{sec:decision} and the fixed-threshold storage gate in \Cref{sec:conformal-gate}. Stage~7, output, executes the implied user-facing action, returning the stored or generated answer or the calibrated refusal. Stage~8 runs at cycle boundaries and consists of the retroactive re-verification pass, the deferred-entry reconsideration pass, the memory consolidation pass, and the regularised fine-tune under the retention guard; the cycle-boundary passes are specified in \Cref{sec:cycle-boundary}.

\subsection{The cycle loop walk-through}
\label{sec:overview-cycle}

A single cycle traces the same trajectory in every iteration. The cycle ingests a fixed workload of fresh queries from the registered training pool. Each query is served by Stages 1 through 7 of the per-query pipeline, accumulating verified episodes in the episodic memory, deferred entries in the bounded deferral buffer, abstentions in the user-facing log, and discards in the discard log. After the workload is exhausted, the cycle-boundary block runs: the retroactive re-verification pass rescores every stored episode under the current generator and applies the upward-only update rule that promotes entries whose composite has crossed the storage threshold and prunes those that have fallen below it; the deferred-entry reconsideration pass examines every entry in the deferral buffer and either promotes, expires, or requeues the entry; the memory consolidation pass collapses near-duplicate stored episodes into a highest-quality representative within each near-duplicate cluster. The fine-tune then draws a training pool from stored episodes whose composite score exceeds the strict training threshold above the storage threshold, applies the parameter-efficient update through the bounded low-rank-adapter envelope on the frozen backbone, and probes the multi-modal retention probe panel to compute the worst per-probe retention ratio. When the retention ratio falls below the registered floor, the cycle aborts and rolls back the weights; the memory is preserved across the rollback so that the verified episodes accumulated during the failed cycle survive, and only the next per-query phase resumes. The cycle then advances, with the locked architectural configuration applied verbatim and only the calibrated probability composite refit on a fresh per-cycle fold. The rationale for this cycle structure is that it decouples the user-facing per-query path from the slow expensive cycle-boundary updates, while the upward-only retroactive rule keeps the stored pool's quality monotonic across cycles even when the verifier itself is improving.

\subsection{A worked example}
\label{sec:overview-worked-example}

The trajectory of two real entries from the cycle-zero artefact bundle anchors the reader before the formal specification in \Cref{sec:design-spec}. The first is a stored FEVER episode and the second is a discarded TriviaQA question; both are recorded with their full per-signal readings in the released cold-start memory catalogue and the released cycle-zero rescored-evaluation catalogue (the literal on-disk paths are listed in the reproducibility appendix).

Consider the FEVER claim \emph{The starting point of the French Revolution was before 1792}, paired with the instruction to label the claim as supports, refutes, or not enough information. Stage~1 computes the pre-routing confidence from the encoder's short-prefix token-probability aggregate and its encoder-layer convergence ratio; the fused and temperature-scaled scalar comfortably clears the safety floor, so no override is triggered. Stage~2 encodes the claim into a sentence-embedding vector and probes the episodic memory. On the cycle-zero pass with an empty memory there is no neighbour to retrieve, so Stage~3 defaults the query to the retrieval-augmented tier; Stage~3b consults the retrieval-utility classifier, which refines the default to the zero-shot tier because the claim is short, self-contained, and rated by the classifier as not benefiting from passage retrieval; Stage~4 generates the label \emph{supports}; Stage~5 evaluates the nine signals across the four families. The top-passage entailment of the answer against the retrieved evidence is high (above $0.94$); the pooled-mean and atomic entailment are lower because only a fraction of the supporting evidence is in the immediate context window; the semantic consistency across resampled chains is high (above $0.85$); the question-answer relevance score is moderate (around $0.69$). The calibrated probability composite aggregates the per-signal calibrated probabilities through the regularised logistic boost layer and produces a storage score of approximately $0.51$. Stage~6 then applies the four-outcome decision tree: the score clears the deferral threshold but falls just below the storage threshold at the locked operating point, so the entry is committed to the deferred buffer for cycle-boundary reconsideration. After the cycle-boundary retroactive pass under the updated generator the same entry crosses the storage threshold and is promoted into the cold-start memory under entry identifier zero. Stage~7 returns the original answer \emph{supports} to the user, which is the correct label.

Consider next the TriviaQA question \emph{The song ``If I Ruled The World'' comes from which musical?}, which appears in the cycle-zero evaluation fold. Stage~1 produces a pre-routing confidence that does not clear the safety floor, because the encoder's convergence ratio on the question is low. The safety override fires on the score-formula path; Stage~3b then consults the retrieval-utility classifier on the same query and confirms the retrieval-augmented dispatch because the base-model self-knowledge probe reads low and the dense passage substrate carries supporting evidence. Stage~4 retrieves passages from the public corpus and runs a grounded generation; the generator returns an empty answer, signalling an explicit refusal. Stage~5 evaluates the verifier signals on the empty answer: the token-level uncertainty signal is very low (around $0.02$), the top-passage entailment is low (around $0.11$), the atomic entailment is low (around $0.09$), and the question-answer relevance score sits at the chance level (around $0.50$). The calibrated probability composite produces a storage score of approximately $0.07$, well below both the deferral and storage cut-points at the locked operating point. Stage~6 then routes the entry: the storage cut-point and the deferral cut-point are both unmet, and the top-passage entailment is below the registered abstention floor, so the decision tree returns the abstention outcome rather than a silent discard. Stage~7 maps the abstention outcome to the user-facing refusal string registered for the deployment and the raw model prediction is suppressed from the user surface. The two trajectories together exercise the architecture's defining behaviour: the verifier-mediated path that promotes a tentatively answered claim to memory after retroactive review, and the verifier-mediated path that turns a low-confidence question into a calibrated refusal rather than a fabricated answer.

\section{Preliminary Design or Design (Model) Specification}
\label{sec:design-spec}

% The high-level architecture figure preserved from the earlier draft.
\begin{figure}[p]
\centering
\resizebox{1.08\textwidth}{!}{%
\begin{tikzpicture}[
    xshift=-0.9cm,
    font=\normalsize,
  % ---- Node styles ----
  stage/.style  = {rectangle, draw=teal!65!black, line width=0.9pt, rounded corners=5pt,
                   minimum height=1.5cm, minimum width=6.5cm, align=center, fill=teal!8},
  router/.style = {rectangle, draw=teal!65!black, line width=0.9pt, rounded corners=5pt,
                   minimum height=1.8cm, minimum width=8.5cm, align=center, fill=yellow!18},
  mcyl/.style   = {cylinder, draw=teal!65!black, line width=0.9pt,
                   shape border rotate=90, aspect=0.25,
                   minimum height=3.0cm, minimum width=3.2cm, align=center, fill=teal!12},
  tA/.style     = {rectangle, draw=green!55!black, line width=1pt, rounded corners=5pt,
                   minimum height=1.7cm, minimum width=4.0cm, align=center, fill=green!12},
  tB/.style     = {rectangle, draw=red!55!orange, line width=1pt, rounded corners=5pt,
                   minimum height=1.7cm, minimum width=4.0cm, align=center, fill=red!12!orange!10},
  tC/.style     = {rectangle, draw=red!60!black, line width=1pt, rounded corners=5pt,
                   minimum height=1.7cm, minimum width=4.0cm, align=center, fill=red!10},
  verify/.style = {rectangle, draw=teal!65!black, line width=1pt, rounded corners=5pt,
                   minimum height=4.2cm, minimum width=9.5cm, align=center, fill=teal!5},
  decide/.style = {rectangle, draw=teal!65!black, line width=0.9pt, rounded corners=5pt,
                   minimum height=1.8cm, minimum width=9.5cm, align=center, fill=yellow!18},
  oServe/.style = {rectangle, draw=green!60!black, line width=1pt, rounded corners=4pt,
                   minimum height=1.4cm, minimum width=3.6cm, align=center,
                   fill=green!30, font=\small\scshape},
  oStore/.style = {rectangle, draw=green!60!black, line width=0.9pt, rounded corners=3pt,
                   minimum height=1.2cm, minimum width=2.7cm, align=center,
                   fill=green!28, font=\small\scshape},
  oDefer/.style = {rectangle, draw=green!45!black, line width=0.9pt, rounded corners=3pt,
                   minimum height=1.2cm, minimum width=2.7cm, align=center,
                   fill=green!12, font=\small\scshape},
  oAbst/.style  = {rectangle, draw=orange!75!black, line width=0.9pt, rounded corners=3pt,
                   minimum height=1.2cm, minimum width=2.7cm, align=center,
                   fill=orange!28, font=\small\scshape},
  oDisc/.style  = {rectangle, draw=red!60!black, line width=0.9pt, rounded corners=3pt,
                   minimum height=1.2cm, minimum width=2.7cm, align=center,
                   fill=red!28, font=\small\scshape},
  simp/.style   = {rectangle, draw=violet!60!black, line width=1pt, rounded corners=5pt,
                   minimum height=3.2cm, minimum width=17.0cm, align=center, fill=violet!7},
  % ---- Arrow styles ----
  farr/.style   = {-Latex, line width=0.9pt},
  garr/.style   = {-Latex, line width=1.1pt, green!50!black},
  oarr/.style   = {-Latex, line width=1.1pt, color=red!55!orange},
  rarr/.style   = {-Latex, line width=1.1pt, red!60!black},
  darr/.style   = {<->, dashed, line width=0.9pt, teal!55!black},
  varr/.style   = {-Latex, line width=1pt, dashed, violet!60!black},
    elab/.style   = {font=\small\itshape, inner sep=2pt, fill=none, draw=none}
]

% ==================================================================
% NODES   (main trunk centred at x = 6; y compressed to fit A4 page)
% ==================================================================

\node[stage]  (q)   at (6,   0.0)  {\textbf{Query $q$}};

\node[stage]  (s1)  at (6,  -2.7)  {%
    \textbf{S1\quad Pre-route confidence}\\[4pt]
    $u_{\text{pre}} = 0.60\,u_{\text{token}} + 0.40\,(1{+}C_{\text{conv}})^{-1}$\\
    {\scriptsize two forward passes;\;$T$ re-fit per cycle}};

\node[stage]  (s2)  at (6,  -5.8)  {%
    \textbf{S2\quad Episodic memory retrieval}\\[4pt]
    SBERT 768-d cosine nearest-neighbour search};

% Episodic memory cylinder  (right of S2, connected by dashed bidirectional arrow)
\node[mcyl]   (mem) at (16.5,-5.8) {%
    \textbf{Episodic}\\[2pt]\textbf{Memory}\\[3pt]$K{=}10^{6}$};

\node[router] (s3)  at (6, -9.2)  {%
    \textbf{S3\quad Three-tier router}\\[4pt]
    $S(q)=0.70\cdot\mathrm{sim}(q,e^{*})+0.30\cdot u_{\text{stored}}(e^{*})$\\
    {\scriptsize safety override:\;if $u_{\text{pre}}<0.60$ the score is bypassed and Tier~3 is forced}};

% ---- Tier row  (T1 far-left / T2 centre / T3 far-right)
%  S5 min-width = 9.5 cm centred at x=6  ->  left edge  ~ x = 1.25
%  T1 placed at x=0 is therefore outside S5, allowing a clean bypass path.
\node[tA] (t1) at (0.0, -12.2)  {%
    \textbf{Tier 1}\\[3pt]return $e^{*}$ answer\\{\scriptsize(from verified memory)}};
\node[tB] (t2) at (6.0, -12.2)  {%
    \textbf{Tier 2}\\[3pt]one-shot CoT generation\\{\scriptsize(base model only)}};
\node[tC] (t3) at (12.0,-12.2)  {%
    \textbf{Tier 3}\\[3pt]RAG: DPR top-5\\{\scriptsize$\to$ cross-encoder rerank to top-3}};

% ---- S4: Tier execution grouping (compact left-side label; no wrapper) ----
\node[draw=none, inner sep=6pt, fit=(t1)(t2)(t3),
      label={[font=\small\bfseries, teal!65!black, anchor=east, xshift=-6pt]left:S4}]
     (s4box) {};

% Tier 1 direct-serve terminal  (left side, same y-band as S5)
\node[oServe] (serve) at (-3.5,-17.0) {%
    \textsc{Answer served}\\{\scriptsize from verified memory}};

% S5 centred at x=6, min-width=9.5 cm  ->  spans ~ x=1.25..10.75
\node[verify] (s5)  at (6, -17.0)  {%
    \textbf{S5\quad Nine-signal Verifier}\\[6pt]
    \textit{Internal:}\quad$u_{\text{token}},\;u_{\text{dropout}},\;u_{\text{internal}}$\\[3pt]
    \textit{Sample-set:}\quad$s_{\text{avg}},\;p_{\text{entail}},\;h_{\text{norm}}$\\[3pt]
    \textit{Grounding:}\quad$p_{\text{ground}}^{\max},\;p_{\text{ground}}^{\text{mean}},\;p_{\text{ground}}^{\text{atomic}}$\\[5pt]
    {\footnotesize\textbf{confabulation gate:}\;
     $u_{\text{internal}}{\geq}0.70\;\wedge\;p_{\text{ground}}^{\max}{\leq}0.20
     \;\Rightarrow\;$\textsc{Discard (early exit)}}};

\node[decide] (s6)  at (6, -21.0)  {%
    \textbf{S6\quad Four-outcome storage decision}\\[5pt]
    {\scriptsize
    $u_{\text{stored}}{\geq}0.65\Rightarrow\textsc{Store}$\quad
    $u_{\text{stored}}{\geq}0.45\Rightarrow\textsc{Deferred}$\quad
    $p_{\text{ground}}^{\max}{<}0.20\Rightarrow\textsc{Abstain}$\quad
    \textit{else:}\;\textsc{Discard}}};

% Terminal outcome boxes  (y = -24.5)
\node[oStore] (out1) at (2.0, -24.5)  {\textsc{Store}\\{\tiny write + train-flag}};
\node[oDefer] (out2) at (5.0, -24.5)  {\textsc{Deferred}\\{\tiny write, no train}};
\node[oAbst]  (out3) at (8.0, -24.5)  {\textsc{Abstain}\\{\tiny principled refusal}};
\node[oDisc]  (out4) at (11.0,-24.5)  {\textsc{Discard}\\{\tiny suppressed}};

% ---- S7: Memory write grouping (compact left-side label; no wrapper) ----
\node[draw=none, inner sep=6pt, fit=(out1)(out2)(out3)(out4),
      label={[font=\small\bfseries, teal!65!black, anchor=east, xshift=-6pt]left:S7}]
     (s7box) {};

% S8: between-cycle self-improvement  (wide band at bottom)
\node[simp] (s8) at (6, -27.5) {%
    \textbf{S8\quad Between-cycle self-improvement loop}
    \quad\textit{(runs once at each cycle boundary, not per query)}\\[6pt]
    (1)~Retroverify all stored entries
        ($\tau_{\text{retro}}{=}0.50$;\;upward-only: never demotes a previously verified entry)\\[3pt]
    (2)~Curate training pool ($\tau_{\text{train}}{=}0.75$;\;90:10 task-to-general mix)\quad
    (3)~Fine-tune:\;$\mathcal{L}(\theta)=\mathcal{L}_{\text{CE}}
        +\tfrac{\lambda}{2}\|\theta-\theta_{\text{prev}}\|^2$,\;$\lambda{=}0.01$\\[3pt]
    (4)~Abort guard:\;$\rho=A_{\text{post}}/A_{\text{pre}}\geq0.93$;\;
        roll back to $\theta_{\text{prev}}$ if violated};

% ==================================================================
% ARROWS
% ==================================================================

% ---- Main trunk ----
\draw[farr] (q.south)  -- (s1.north);
\draw[farr] (s1.south) -- node[elab,right]{$u_{\text{pre}}$} (s2.north);
\draw[farr] (s2.south) -- node[elab,right]{$u_{\text{pre}}$ + query embedding} (s3.north);

% ---- S2 <-> Episodic memory  (bidirectional dashed) ----
\draw[darr] (s2.east)
    -- node[elab,above,pos=0.55,align=center]{%
        retrieve $e^{*}$\;:\;$\mathrm{sim}(q,e^{*})$,\\$u_{\text{stored}}(e^{*})$}
    (mem.west);

% ---- S3 -> three tiers ----
\coordinate (br) at (6, -10.9);   % branch waypoint below S3
\draw[farr] (s3.south) -- (br);

% Tier 1 (green): branch LEFT, condition label above horizontal leg
\draw[garr]
    (br) --
    node[elab,above,pos=0.45]{\textbf{if}\;$S(q)\geq0.90$}
    (0.0,-10.9) -- (t1.north);

% Tier 2 (orange): straight down, condition label to the RIGHT
\draw[oarr]
    (br) --
    node[elab,right,pos=0.5]{\textbf{elif}\;$\mathrm{sim}(q,e^{*})\geq0.75$}
    (t2.north);

% Tier 3 (red): branch RIGHT, condition label above horizontal leg
\draw[rarr]
    (br) --
    node[elab,above,pos=0.75]{\textbf{else};\;or forced when\;$u_{\text{pre}}<0.60$}
    (12.0,-10.9) -- (t3.north);

% ---- Tier 1 bypass: route LEFT then DOWN, well outside S5 ----
% S5 left edge ≈ x=1.25; this path uses x=-3.5 to stay clear.
\coordinate (t1byp1) at (0.0, -14.0);    % just below T1
\coordinate (t1byp2) at (-3.5,-14.0);    % far left
\draw[garr]
    (t1.south) -- (t1byp1) --
    node[elab,below,pos=0.5]{\textit{bypass S5 and S6: episode already verified at storage time}}
    (t1byp2) -- (serve.north);

% ---- T2 -> S5 ----
\draw[oarr]
    (t2.south) --
    node[elab,right,pos=0.45]{answer + reasoning chain}
    (t2 |- s5.north);

% ---- T3 -> S5 routed on the right to avoid center-label overlap ----
\draw[rarr]
    (t3.south) --
    (12.0,-14.2) --
    (13.6,-14.2) --
    node[elab,right,pos=0.55]{answer + reasoning chain}
    (13.6,-16.9) --
    (s5.east);

% ---- S5 -> S6 ----
\draw[farr]
    (s5.south) --
    node[elab,right]{outputs:\;$u_{\text{stored}},\;p_{\text{contra}},\;p_{\text{ground}}^{\max}$}
    (s6.north);

% ---- S6 fan-out -> terminal boxes ----
\coordinate (br2) at (6,-22.8);
\draw[farr] (s6.south) -- (br2);
\draw[farr] (br2) -| (out1.north);
\draw[farr] (br2) -| (out2.north);
\draw[farr] (br2) -| (out3.north);
\draw[farr] (br2) -| (out4.north);

% ---- Store + Deferred accumulate into S8 at cycle end ----
% Drop only 0.3 cm so the horizontal segment stays ABOVE s8.north,
% then the vertical leg enters s8 cleanly from the top.
\draw[varr]
    (out1.south) -- ++(0,-0.3) -|
    node[elab,left,pos=0.05]{\textit{cycle end}}
    ($(s8.north)+(-2.5,0)$);
\draw[varr]
    (out2.south) -- ++(0,-0.3) -|
    ($(s8.north)+(-0.8,0)$);

% ---- S8 -> S1 : OUTER TEN-CYCLE LOOP ----
%
%  After S8 concludes, the updated weights θ_new (or θ_prev on rollback)
%  are loaded and Stage 1 of Cycle n+1 begins.  This arrow is the outer loop.
%
%  Route: S8.west  →  far left (x ≈ -5.5)  →  up to y of S1  →  S1.west
%
\draw[varr]
    (s8.west)
    -- ++(-5.5, 0) coordinate (ll)
    -- node[elab,left,pos=0.45,align=left]{%
        \textbf{Cycle $n{+}1$ begins}\\[3pt]
        S1 receives\\
        $\theta_{\text{new}}$ (fine-tuned)\\[1pt]
        or $\theta_{\text{prev}}$ (restored\\
        if abort guard fired)}
    (ll |- s1.west)
    -- (s1.west);

\end{tikzpicture}%
}
\caption{%


\subsection{Generator and retrieval substrate}
\label{sec:gen-retrieval}

The architecture is built on four substrate components whose architectural identity is fixed across every cycle of the self-improvement loop. The generator is a three-billion-parameter open-weight instruction-tuned decoder-only model from the Qwen 2.5 family \cite{qwen25}, used for all generation in the zero-shot tier and the retrieval-augmented tier and used as the source of internal-calibration signals in the verifier. The choice of a small generator rather than a frontier-scale model is deliberate. A small generator admits parameter-efficient fine-tuning on a single consumer-grade graphics processor through low-rank adaptation on the attention and feed-forward projections, with the trainable footprint at roughly two percent of the full backbone parameter count, which is the necessary condition for the cycle-boundary fine-tune to run inside the project's compute envelope. Equally, a small generator's baseline accuracy on the five-benchmark panel registered for the main run is high enough that the verifier's expected false-positive band is non-degenerate, so the storage gate has a usable signal-to-noise ratio without depending on a frontier-scale capacity that the project cannot afford. The adapter parameters are the single locus of cycle-to-cycle drift; the underlying backbone is frozen and everything else in the substrate is held constant.

The dense passage index is a single sentence-embedding-encoded snapshot of a public Wikipedia-derived corpus, structured as a dense-retrieval index in the spirit of \cite{karpukhin-etal-2020-dense}. The encoder is a sentence-transformer model of moderate dimension \cite{reimers-gurevych-2019-sentence} that is also used to encode the user's query at lookup time, so query and passage embeddings live in the same space. The index is a tiered structure: a flat in-memory index for capacity below the registered switching threshold, promoted to an inverted-file product-quantised index above that threshold so that lookup latency stays within the per-query envelope at the deployment-scale memory size projected for production. The retrieval substrate is shared by the retrieval-augmented tier and by the external grounding family of verifier signals, so a single passage retrieval per query feeds both the grounded generation and the entailment scoring.

The natural-language-inference judge that the verifier consumes for chain-to-answer entailment and for the external grounding signals is realised by the MiniCheck checkpoint \cite{tang2024minicheck}, a unary supported-against-context scorer trained on language-model-generated claim-support data and therefore distributionally matched to the kind of generated answers the verifier is asked to score. The judge is a unary model that returns a single supported probability rather than a three-way label distribution, which is why the contradiction signal of the verifier ensemble has been retired at the architectural level: the contradiction class never receives non-zero mass under this judge.

The cross-encoder reranker is a separately trained encoder-only model from the bge-reranker family that scores text pairs by relevance. It is used in two places in the architecture, both as a single shared instance to amortise its load cost: as the optional reranker on the retrieval-augmented tier and on the external grounding family, where the verifier reranks the top-twenty retrieved passages down to the top-three before scoring entailment, and as the question-answer relevance signal in the verifier. The reranker and the natural-language-inference judge are distinct models with distinct training distributions; combining them as separate substrate components is what allows the verifier to keep the question-relevance signal architecturally orthogonal to the grounding family, so that an off-topic answer that is incidentally entailed by a retrieved passage is still caught by the relevance signal. The capacity and latency profile of the substrate as a whole is reported empirically in \Cref{ch:evaluation}, and the rationale for the substrate choices over alternatives is registered against the constraints in \Cref{sec:constraints}: the single-backbone commitment, the local-execution commitment, and the bounded-cycle commitment together fix the substrate to the components specified here.

\subsection{Pre-routing confidence}
\label{sec:pre-routing}

The pre-routing confidence is a single scalar computed from a single short forward pass on the query, used at two points in the per-query pipeline: as a one-dimensional input to the dispatch score that selects the tier, and as the sole input to the safety override that forces the retrieval-augmented tier when the confidence falls below a fixed floor. The signal is a fusion of two cheap quantities, both produced by the same forward pass over the query and a short greedy decode that follows it. The first is a short-greedy-decode token-probability aggregate, computed by running a thirty-two-token greedy decode under the chat template on the query and recording the geometric mean of the per-token log-probabilities of the chosen tokens; this captures how confidently the generator commits to its next-token choices under the question's framing. The second is an encoder-layer convergence ratio, computed from the hidden-state stack of the same forward pass as the variance ratio between the early-layer half and the late-layer half of the decoder; this captures how stably the representation has settled across the layers that drive the downstream confidence, with high variance in the early layers relative to the late layers reading as the model not having converged on a single interpretation of the question. The fusion uses fixed mixing weights of $0.60$ on the token-probability aggregate and $0.40$ on the inverse-one-plus convergence ratio, registered before the calibration phase rather than learned, so the fused scalar's interpretation does not drift with the fine-tune.

The fused scalar is calibrated to a probability through temperature scaling \cite{pmlr-v70-guo17a} fitted on a disjoint calibration fold by minimising binary negative-log-likelihood under the labels-correct-under-zero-shot binary signal. Temperature scaling is preferred over Platt scaling, isotonic regression, or histogram binning at this site for three reasons. Temperature scaling has a single scalar parameter, which prevents the calibration fold from over-fitting at this site and conserves the fold's degrees of freedom for the heavier post-generation calibration work. Temperature scaling preserves the relative ordering of inputs, which matters for the dispatch score that uses the scalar as a continuous input rather than a hard label. Temperature scaling has a closed-form interpretation as a uniform sharpening or softening of the score distribution, which makes the calibrated output diagnosable at each cycle boundary by comparing the temperature scalar across cycles.

The non-trivial design choice is that the pre-routing confidence does not require the multi-chain resampling that the post-generation verifier uses. The cost of the signal is therefore paid per query regardless of which tier eventually serves the answer (a single forward pass over the query plus a thirty-two-token greedy decode), in exchange for the dispatch score and the safety override having a usable confidence signal at every query. The cost of the more expensive sample-based signals, including the dropout-disagreement signal and the sample-set agreement family, is deferred to the post-generation verifier on the zero-shot and retrieval-augmented paths only, so memory-direct queries incur no decode-phase signal cost beyond the short pre-routing decode. The contribution to the hypotheses registered in \Cref{ch:introduction} is direct: the safety override that this signal alone drives is what closes the high-uncertainty failure mode where neither memory nor zero-shot generation has earned the right to answer, and the dispatch score that uses this signal jointly with the memory similarity score is what permits the memory-direct tier to be cheap without sacrificing safety. \Cref{alg:pre-routing} states the fusion procedure as it is implemented in the per-query pipeline.

\begin{algorithm}[!htbp]
\caption{Pre-routing confidence fusion (Stage~1).}
\label{alg:pre-routing}
\begin{algorithmic}[1]
\Require query $q$, generator $G$, max generation length $L = 32$, per-benchmark temperature scalar $T_{b}$, mixing weights $w_{\text{tok}} = 0.60$ and $w_{\text{conv}} = 0.40$
\State $\big(\{t_{i}\}_{i=1}^{n}, \{\log p_{i}\}_{i=1}^{n}\big) \gets G.\mathrm{generate}(q;\ L,\ \text{greedy},\ \text{scores})$ \Comment{short greedy decode of up to $L$ tokens; greedy disables sampling and \emph{scores} returns token log-probabilities}
\State $u_{\text{tok}} \gets \exp\!\left(\tfrac{1}{n}\sum_{i=1}^{n} \log p_{i}\right)$ \Comment{geometric mean of per-token log-probabilities of the chosen tokens}
\State $\mathbf{H} \gets G.\mathrm{forward}(q,\ \texttt{output\_hidden\_states}{=}\text{True})$ \Comment{full hidden-state stack from a single forward pass}
\State $c_{\text{conv}} \gets \mathrm{Var}\big(\mathbf{H}^{(1:L/2)}\big) \;/\; \mathrm{Var}\big(\mathbf{H}^{(L/2+1:L)}\big)$ \Comment{variance ratio of early to late decoder layers, the Nandakishor convergence ratio}
\State $z \gets w_{\text{tok}}\cdot u_{\text{tok}} + w_{\text{conv}}\cdot \tfrac{1}{1+c_{\text{conv}}}$ \Comment{fixed-weight fusion}
\State $u_{\text{pre}} \gets \sigma\!\big(\mathrm{logit}(z)/T_{b}\big)$ \Comment{per-benchmark temperature scaling on the logit}
\State \Return $u_{\text{pre}}$
\end{algorithmic}
\end{algorithm}

\subsection{Episodic memory store}
\label{sec:memory-store}

The episodic memory is the architecture's single source of long-horizon learning between cycle-boundary fine-tunes. The design is in the lineage of memory-augmented neural networks for sequence and question-answering tasks \cite{weston2015memory,graves2014neural,graves2016differentiable}, of writeable episodic key-value stores for language models \cite{larimar2024das}, and of insertion-time memory augmentation \cite{xu2025amem}, but is specialised here as an external, query-keyed episodic store whose distinguishing property is the scalar storage score persisted with each entry. Each entry is an indivisible record that pairs the question, the served answer, the supporting reasoning chain, the sentence-embedding vector of the question, the verifier's per-signal readings at storage time, the calibrated storage score, and the entry's storage cycle. Storing the question and the answer together with the supporting context and the per-signal readings allows the cycle-boundary retroactive re-verification pass to recompute the storage score under the updated generator without re-running the original generation, which is the architectural property that makes upward-only updates affordable at scale.

Similarity search is performed by cosine similarity in the sentence-embedding space, against a tiered index that promotes from a flat in-memory structure to an inverted-file product-quantised structure above a registered switching threshold on the entry count. Cosine similarity is preferred over Euclidean distance because the sentence-transformer encoder is trained under a contrastive objective that places semantically equivalent questions on the same direction in embedding space, so angular separation tracks semantic distance more reliably than $L_2$ separation. The tiered index choice is justified by the deployment-scale projection: at the workloads reported in this thesis the memory grows to the low tens of thousands of entries and the flat index suffices, but the deployed architecture has a path to the millions-of-entries regime under the inverted-file quantised index without a code rewrite.

A novelty filter at write time bounds the rate at which the memory absorbs near-duplicate entries. When a candidate write's nearest neighbour in the existing memory exceeds the novelty similarity threshold (set to a value above which two entries are considered paraphrases of the same query), the candidate is suppressed; this prevents the memory from bloating with paraphrase-level duplicates that consume capacity without adding coverage. The threshold is registered before the calibration phase and held constant across cycles. A complementary consolidation pass runs at every cycle boundary at a more conservative similarity threshold, which collapses near-duplicate clusters by retaining the highest-quality representative within each cluster and merging the source-benchmark provenance into the surviving entry.

A capacity policy bounds the memory at a registered upper limit on the number of stored entries. When the limit is reached, a value-score pruning rule selects entries for removal: the value score combines retrieval frequency, success rate at retrieval time, recency, and storage-time confidence into a single rank, and the lowest-ranked entries are pruned until the memory is back below the limit. The pruning rule is invoked only when the limit is exceeded, so the small memories reported in this thesis never trigger it; the policy nonetheless governs the architecture's deployment-scale behaviour and is verified empirically through the consolidation diagnostics in \Cref{ch:evaluation}.

The contribution of the memory design to the registered hypotheses is twofold. First, by storing the per-signal readings together with the answer, the memory supports the upward-only retroactive update rule that gives the pool-purity lower bound of \Cref{thm:purity} its monotonic structure across cycles. Second, by enforcing the novelty and consolidation policies, the memory bounds the rate at which storage capacity accumulates without coverage, which is what makes the deployment-cost projection in \Cref{sec:economic} non-trivial: the per-query cost saving depends on memory entries being distinct from one another at retrieval time.

\subsection{Three-tier router with the safety override}
\label{sec:router}

The router is the dispatch logic that sits between the encode-and-lookup stage and the answer-generating tiers. It takes three inputs from the upstream stages: the cosine similarity of the query to the nearest stored entry, the stored-quality score of that neighbour, and the calibrated pre-routing confidence; it produces one output, the tier index, together with a flag that records whether the safety override fired.

The combined dispatch score blends the similarity score and the neighbour's stored-quality score under a fixed weight registered before the calibration phase. The combined score uses a similarity weight of $\lambda = 0.70$ on the cosine similarity and the complementary weight $1-\lambda = 0.30$ on the neighbour's stored-quality score. The choice to give similarity the larger weight reflects the empirical fact that retrieval errors at the lookup step degrade Tier 1 quality more than uncertainty in the stored-quality score does, because a wrong neighbour is a categorically different query while an under-confident correct neighbour is still a correct answer; the rationale for the specific weight is registered in the calibration runbook and the sensitivity analysis is reported in \Cref{ch:evaluation}.

The dispatch rules execute in a fixed order. The safety override is checked first: when the pre-routing confidence falls below the registered floor of $0.38$, the router dispatches to the retrieval-augmented tier and records the override flag, regardless of the combined score. The floor was lowered from the originally-registered value of $0.60$ during the cycle-1 routing audit, which observed that the higher floor was forcing well-supported queries onto the retrieval-augmented tier under the post-temperature-cap pre-routing distribution and depressing the architecture's effective tier utilisation; the lowered floor preserves the architectural intent of the safety override while admitting a wider band of queries to the cheaper tiers. The ordering matters because the override expresses an architectural constraint that the system must not generate without grounded retrieval whenever its pre-generation readiness is low. If the override does not fire, the combined score is compared against the Tier 1 threshold of $0.90$; queries above the threshold are dispatched to the memory-direct tier, which serves the neighbour's stored answer without invoking the generator. Queries that do not clear the Tier 1 threshold are checked against the Tier 2 similarity threshold of $0.75$ on the similarity score alone; queries above this threshold are dispatched to the zero-shot tier, which runs a generation on the fine-tuned generator without grounded retrieval. Queries that do not clear either threshold fall through to the retrieval-augmented tier, which retrieves passages from the public corpus and runs a grounded generation. The retrieval-augmented tier is the architectural floor for the routing rule of this subsection: it is the path that the system falls back to whenever the cheaper paths cannot earn the right to answer, subject to the per-query retrieval-utility refinement of \Cref{sec:ruc} that may redirect the safety-veto fall-through and the score-formula fall-through to the zero-shot tier whenever retrieval is judged unlikely to help.

The three-mechanism separation between the combined-score dispatch, the safety override, and the per-query retrieval-utility decision of \Cref{sec:ruc} is what distinguishes the router from a single-threshold dispatcher and from prior three-tier routers such as Adaptive-RAG \cite{jeong2024adaptive}, whose three tiers are gated by a query-complexity classifier rather than by per-entry stored quality on a shared memory. A single-threshold dispatcher would conflate the question of whether the memory has a useful neighbour with the question of whether the model is ready to generate, and would either over-retrieve under high memory similarity with low pre-routing confidence or under-retrieve in the opposite regime. The three-mechanism design keeps the three questions separate: the combined score captures the memory's coverage of the incoming query, the safety floor captures the model's readiness to generate, the retrieval-utility decision captures whether retrieval is expected to help on the particular query, and a query is dispatched to the cheap memory-direct tier only when the first two mechanisms agree that the cheap path is safe while the third mechanism governs the choice between zero-shot and retrieval-augmented generation on the residual queries.

The expected per-tier hit rate envelope predicted by the architecture is read off the dispatch rules together with the self-improvement loop's effect on the generator's zero-shot capability. Early in a cycle, when the memory is small and the adapter has not yet absorbed cycle-zero verified reasoning, the Tier 1 share is bounded above by the fraction of queries with a stored neighbour whose combined score clears $0.90$; the Tier 2 share is bounded above by the fraction of queries with a similar neighbour clearing $0.75$ but no Tier 1-quality match, plus the share of residual queries that the retrieval-utility classifier of \Cref{sec:ruc} redirects from the retrieval-augmented fallback to the zero-shot tier; the Tier 3 share absorbs the remaining residual on which the classifier judges retrieval helpful. As cycles progress, two coupled redistributions reduce the Tier 3 share. The Tier 1 share grows because the verified episodic store accumulates cached answers for recurring queries that earlier cycles served through Tier 3 retrieval-augmented generation, and the retroactive re-verification pass keeps the cached composites fresh under the updated generator. The Tier 2 share grows because the self-improvement loop folds verified retrieval-augmented reasoning into the generator's parametric capacity at every cycle boundary, so queries that would have required Tier 3 retrieval at cycle zero are answered zero-shot from the fine-tuned generator at later cycles without paying the retrieval cost. The empirical realisation of this envelope is reported in \Cref{ch:evaluation} as the per-tier hit rate trajectory. \Cref{tab:dispatch-rules} summarises the dispatch rules in the order of evaluation, and \Cref{alg:router} states the routing procedure as implemented.

\begin{table}[!htbp]
\centering
\small
\begin{tabular}{@{}p{0.6cm}p{1.0cm}p{6.5cm}p{4.5cm}@{}}
\toprule
\textbf{Order} & \textbf{Tier} & \textbf{Condition} & \textbf{Output path} \\
\midrule
1 & Tier 3 & $u_{\text{pre}} < 0.38$ (safety override) & Retrieval-augmented generation \\
2 & Tier 1 & $\mathcal{S} = 0.7\,s + 0.3\,u_{\text{stored}} \ge 0.90$ & Serve neighbour's stored answer \\
3 & Tier 2 & $s \ge 0.75$ & Zero-shot generation \\
4 & Tier 3 & otherwise & Retrieval-augmented generation (fallback) \\
\bottomrule
\end{tabular}
\caption{Three-tier dispatch rules in order of evaluation. The safety override is checked first; the combined-score Tier 1 rule is checked second; the similarity-only Tier 2 rule is checked third; the retrieval-augmented tier absorbs the remainder, subject to the per-query retrieval-utility refinement of \Cref{tab:ruc-dispatch}.}
\label{tab:dispatch-rules}
\end{table}

\begin{algorithm}[!htbp]
\caption{Three-tier router with safety override (Stage~3).}
\label{alg:router}
\begin{algorithmic}[1]
\Require similarity $s$, neighbour stored-quality $u_{\text{stored}}$, pre-routing confidence $u_{\text{pre}}$
\Require constants $\lambda = 0.70$, $\theta_{1} = 0.90$, $\theta_{\text{sim}} = 0.75$, $\phi_{\text{safe}} = 0.38$
\State $\mathcal{S} \gets \lambda\, s + (1-\lambda)\, u_{\text{stored}}$ \Comment{combined dispatch score}
\If{$u_{\text{pre}} < \phi_{\text{safe}}$}
    \State \Return $(\text{Tier 3},\ \texttt{safety\_override} = \text{True})$ \Comment{refined downstream by \Cref{alg:ruc}}
\ElsIf{$\mathcal{S} \ge \theta_{1}$}
    \State \Return $(\text{Tier 1},\ \texttt{safety\_override} = \text{False})$
\ElsIf{$s \ge \theta_{\text{sim}}$}
    \State \Return $(\text{Tier 2},\ \texttt{safety\_override} = \text{False})$
\Else
    \State \Return $(\text{Tier 3},\ \texttt{safety\_override} = \text{False})$ \Comment{refined downstream by \Cref{alg:ruc}}
\EndIf
\end{algorithmic}
\end{algorithm}

\subsection{Retrieval-utility classifier}
\label{sec:ruc}

The router of \Cref{sec:router} dispatches every query that fails the Tier 1 combined-score gate and the Tier 2 similarity gate to the retrieval-augmented tier as the architectural fallback. The dispatch is the correct default whenever the dense passage substrate carries grounding evidence that helps the generator answer the query; it is the wrong default whenever the retrieved evidence amplifies the question's misconception framing or distracts the generator from a stronger parametric answer. The empirical anchor for the wrong-default mode is the regression on misconception-bait benchmarks observed in the pre-rerun trajectory of the architecture without a routing refinement: TruthfulQA exact match under the Haiku-judged correctness rule declined cycle over cycle as the self-improvement loop folded retrieval-amplified answers into the verified pool. The same architectural diagnosis was reached independently by the selective-retrieval line of work reviewed in \Cref{sec:lit-selective-retrieval}: Mallen et al.~\cite{mallen2023whennot} on PopQA, Wang et al.~\cite{wang2023skr} on self-knowledge routing, Jeong et al.~\cite{jeong2024adaptive} on Adaptive-RAG, and Baek et al.~\cite{baek2025probing} on Probing-RAG all converge on the architectural conclusion that an unconditional retrieval dispatch leaves measurable accuracy on the table relative to a per-query routing decision. The retrieval-utility classifier is CAEM's architectural correction in that lineage. It sits between the two-mechanism dispatcher of \Cref{sec:router} and the retrieval-augmented fallback, reads a per-query feature vector built from the query, the top-ranked passage, and a base-model self-knowledge probe, and decides whether the query is one on which retrieval helps the generator or one on which retrieval hurts. On the former it forwards the query to the retrieval-augmented tier as the router intended; on the latter it routes the query to the zero-shot tier, so that the generator's parametric knowledge serves the answer without paying the retrieval cost and without admitting passage-amplified errors into the self-improvement loop.

The classifier is a binary head with the two output classes \textsc{rag} (route to the retrieval-augmented tier) and \textsc{direct} (route to the zero-shot tier). It produces a calibration score $p_{\text{RAG}} \in [0, 1]$ from a thirty-six-dimensional feature vector through standard-scaler normalisation followed by a logistic-regression scoring head with L2 regularisation. The deployed routing rule compares $p_{\text{RAG}}$ against a registered threshold $\tau_{\text{RUC}} = 0.375$ and dispatches to the retrieval-augmented tier when $p_{\text{RAG}} \ge \tau_{\text{RUC}}$ and to the zero-shot tier otherwise. The threshold is well below the unbiased operating point of $0.5$ because the cost of a missed retrieval-helpful query (a query that needed retrieval but was sent to zero-shot generation) is asymmetric with the cost of a missed retrieval-hurt query (a query that should have stayed on zero-shot but was sent to retrieval-augmented generation): the first failure mode under-utilises retrieval evidence that was available, while the second admits passage amplification into a hard distribution where it had been measured to hurt. The threshold sweep registered at training time selects the operating point that jointly maximises the count of correctly-routed retrieval-helpful queries and the count of correctly-routed retrieval-hurt queries on the held-out folds; the realised maximum sits at the asymmetric value above.

The thirty-six features partition into three groups. The first group is computed inline from the question text and is always available at routing time: the question token length, a one-hot encoding over the seven canonical interrogative templates (\emph{what}, \emph{when}, \emph{where}, \emph{who}, \emph{why}, \emph{how}, \emph{yes-or-no}), the presence of negation and of temporal or numerical cues, an entity count, a myth-regex hit indicator for the canonical misconception phrasings, and three linguistic markers (\emph{adversarial phrasing}, \emph{opinion-seeking}, \emph{open-ended}) drawn from the question's surface form. The second group is computed from the top-ranked passage that the dense retriever returned for the query: the top-one similarity score, the top-one passage--question entity overlap, three answer-type matches between the question's expected answer type and the passage's surface (year-match, number-match, and a numeric-presence flag), and a Wikidata-entity-presence indicator that records whether the question's primary entity appears in the structured knowledge base. The third group is a battery of retrieval-quality and self-knowledge features that the routing-time pipeline assembles from the same retrieval substrate the retrieval-augmented tier would use: the top-five FAISS similarity statistics (maximum, mean, standard deviation), the cross-encoder rerank statistics at the same top-five (top-one rerank score and the mean and standard deviation over the top-five), three pairwise natural-language-inference statistics computed on the ten pairs of the top-five passages (mean, minimum, and standard deviation of the entailment probability), and a Kadavath-style self-knowledge probe $p_{\text{IK}}$ that reads the base generator's last-token hidden state on the question alone, without adapter, and predicts whether the model knows the answer parametrically. The third group is the load-bearing axis of the classifier: the probe $p_{\text{IK}}$~\cite{kadavath2022language} carries the largest absolute coefficient in the trained scoring head ($-1.77$), so high parametric self-knowledge pushes the classifier strongly toward \textsc{direct}, while passage-quality features with positive coefficients push toward \textsc{rag} only when the parametric self-knowledge is low.

The training set is constructed from a content-hash-disjoint fresh pool that does not intersect the calibration fold, the seed pool, the per-cycle stream chunks, or the evaluation and test folds the rest of the architecture consumes. Each candidate row in the fresh pool is scored under the two routing destinations the classifier decides between: the question is answered both by a Tier 2 baseline (a zero-shot generation against the same backbone, without retrieval) and by a Tier 3 baseline (a retrieval-augmented generation against the same backbone, with the same dense substrate the deployment uses), and the two exact-match scores are recorded. The row's label is then assigned by a deterministic rule on the two exact-match scores: the row is a positive (\emph{retrieval-helpful}, label one) when the Tier 3 exact match exceeds the Tier 2 exact match, a negative (\emph{retrieval-hurt}, label zero) when Tier 2 exact match exceeds Tier 3, and is dropped when the two scores tie. The drop rule discards both kinds of ties: rows where both routings produced a correct answer (the classifier's decision is irrelevant) and rows where both routings produced a wrong answer (the classifier's decision is also irrelevant because the architecture's downstream gate will reject both at the storage step). The architectural intent of the drop rule is to give the classifier a training signal that strictly disambiguates the cases where the routing decision matters, and to let the architecture's downstream gate absorb the cases where it does not. The realised dataset spans eight benchmarks (FEVER, TriviaQA, Natural Questions, HaluEval-QA, CommonsenseQA, OpenBookQA, TruthfulQA, and StrategyQA), with three thousand two hundred and eight rows after the drop rule, of which one thousand eight hundred and forty-six are retrieval-helpful and one thousand three hundred and sixty-two are retrieval-hurt, a ratio of approximately $1.36$ to one. TruthfulQA exact match in the label assembly uses the Haiku-judged correctness rule rather than the default ROUGE proxy, consistent with the architecture's downstream TruthfulQA scoring contract.

Three benchmarks in the training set (Natural Questions, HaluEval-QA, OpenBookQA) are not part of the five-benchmark evaluation panel of \Cref{ch:evaluation}: they are present in the classifier's training distribution only to broaden the coverage of retrieval-utility regimes the head must learn. This is a deliberate asymmetry between the classifier's training distribution and the architecture's evaluation distribution: the classifier learns from a broader pool than the architecture is scored on, so that the classifier's per-benchmark generalisation onto the evaluation panel is the load-bearing receipt the empirical chapter audits rather than an in-distribution recall measurement. The three training-only benchmarks do not feed any other component of the architecture, in particular the self-improvement training pool, so the asymmetry is contained within the classifier's training signal and does not leak into any downstream subsystem.

The classifier is fit under a nested five-fold stratified cross-validation. The outer folds partition the three thousand two hundred and eight rows by label into five equally-sized strata, holding out one stratum per outer split for an out-of-fold prediction; the inner folds inside each training stratum select the L2 regularisation strength from the grid $\{0.005, 0.01, 0.05, 0.1, 0.5, 1.0, 5.0\}$ by the same stratified protocol on the inner training data, so that the regularisation choice is made without ever seeing the outer held-out fold. The final-fit regularisation strength is read off the mode of the per-outer-fold inner selections and lands at $C = 0.5$, paired with a class-balancing weight that compensates for the asymmetric label rate. The held-out operating-point threshold $\tau_{\text{RUC}}$ is then selected by a sweep over the candidate grid $\{0.20, 0.225, \ldots, 0.75\}$ on the concatenated out-of-fold predictions, choosing the value that maximises the joint correctly-routed count across the two label classes. The nested protocol decouples the regularisation choice from the threshold choice and prevents either choice from leaking the held-out fold into the trained classifier.

Three held-out scores anchor the classifier's calibration receipt and license the classifier as a deployment-ready architectural component. The pooled out-of-fold AUROC across the five outer folds is $0.850$, with a per-fold standard deviation of $0.008$ across the five outer folds and an in-sample-versus-out-of-fold drift of $0.011$ that closes well inside the registered overfit ceiling. At the deployed threshold the joint correctly-routed accuracy is $77.2\%$ of the three thousand two hundred and eight held-out rows, with a confusion matrix that reads one thousand five hundred and fifty-four correctly-routed retrieval-helpful queries against two hundred and ninety-two miss-routed ones (recall $0.842$ on the retrieval-helpful class) and nine hundred and twenty-three correctly-routed retrieval-hurt queries against four hundred and thirty-nine miss-routed ones (specificity $0.678$ on the retrieval-hurt class). Per-benchmark AUROC on the same out-of-fold predictions ranges from $0.686$ on OpenBookQA at the low end to $0.936$ on TriviaQA at the high end; the realised band on the five-benchmark evaluation panel reads FEVER $0.718$, TriviaQA $0.936$, CommonsenseQA $0.728$, TruthfulQA $0.732$, and StrategyQA $0.852$. The TruthfulQA reading is the load-bearing transfer score because the regression on TruthfulQA is the empirical anchor that motivated the classifier in the first place; the realised $0.732$ sits below the stretch target of $0.75$ that the pre-registered classifier acceptance gate registered, and is reported in \Cref{ch:evaluation} alongside the realised TruthfulQA recovery that the classifier delivers under deployment. The self-knowledge probe $p_{\text{IK}}$ used as a feature carries its own held-out calibration receipt: the probe is fit on eleven thousand nine hundred and eighty questions under the same nested five-fold protocol, on the binary task of predicting whether the base generator answers the question correctly without retrieval, with a mean out-of-fold AUROC of $0.771$.

The classifier consumes a frozen scoring head at deployment. It is fit once at the cycle-zero pre-launch stage on the fresh pool above and is not refit at any cycle boundary across the realised six-cycle trajectory. The frozen choice is deliberate: the classifier's input features are computed from quantities the architecture does not retrain (the question text, the dense passage substrate, the base generator's hidden state under the disable-adapter context, and a pre-trained cross-encoder reranker), so cycle-boundary drift in the classifier's input distribution is bounded by the architecture's substrate-level stability rather than by the adapter-level drift the self-improvement loop induces. The architectural advantage of the frozen design is that the classifier's per-cycle behaviour can be audited as a pure routing-distribution drift on a fixed scoring head rather than as a coupled drift in both the input distribution and the scoring head's coefficients. The corresponding architectural cost is that any benchmark-level distribution shift between the fresh pool and a downstream evaluation pool is absorbed entirely by the classifier's per-benchmark generalisation rather than by a cycle-boundary refit, and the receipt for this absorption is the per-benchmark AUROC band on the evaluation panel reported above and the routing distribution trajectory reported in \Cref{ch:evaluation}. The lightweight-classifier-only operating point CAEM adopts is consistent with the empirical conclusion of Moskvoretskii et al.~\cite{moskvoretskii2025llm}, who show that an external classifier with no language-model forward pass at routing time matches the routing accuracy of language-model-based selective retrievers on the standard suite; CAEM differs in that the architecturally heavy retrieval-quality features (cross-encoder rerank, pairwise NLI on top-five passages) are reused from the verifier ensemble's already-loaded substrate rather than computed afresh.

The router consumes the classifier at exactly two of the four dispatch outcomes of \Cref{tab:dispatch-rules}: the safety-override outcome and the score-formula fall-through outcome. Both outcomes share the property that the pre-classifier dispatch would unconditionally route to the retrieval-augmented tier; the classifier intercepts each one and gives it a chance to be redirected to the zero-shot tier instead. The Tier 1 outcome (the combined-score memory-direct dispatch) and the Tier 2 outcome (the similarity-only zero-shot dispatch) are never gated by the classifier, because both already serve a tier other than the retrieval-augmented fallback and the classifier's binary choice is between exactly those two tiers. The post-classifier dispatch order is summarised in \Cref{tab:ruc-dispatch} and the classifier's scoring and dispatch logic is stated in \Cref{alg:ruc}.

\begin{table}[!htbp]
\centering
\small
\begin{tabular}{@{}p{0.6cm}p{1.0cm}p{6.5cm}p{4.5cm}@{}}
\toprule
\textbf{Order} & \textbf{Tier} & \textbf{Condition} & \textbf{Output path} \\
\midrule
1 & Tier 3 / 2 & $u_{\text{pre}} < 0.38$ (safety override): consult classifier; tier $= 3$ if $p_{\text{RAG}} \ge \tau_{\text{RUC}}$ else tier $= 2$ & Retrieval-augmented or zero-shot generation, gated by the classifier \\
2 & Tier 1 & $\mathcal{S} = 0.7\,s + 0.3\,u_{\text{stored}} \ge 0.90$ & Serve neighbour's stored answer \\
3 & Tier 2 & $s \ge 0.75$ & Zero-shot generation \\
4 & Tier 3 / 2 & otherwise (score-formula fall-through): consult classifier; tier $= 3$ if $p_{\text{RAG}} \ge \tau_{\text{RUC}}$ else tier $= 2$ & Retrieval-augmented or zero-shot generation, gated by the classifier \\
\bottomrule
\end{tabular}
\caption{Post-classifier dispatch rules in order of evaluation. The retrieval-utility classifier intercepts the two dispatch outcomes (safety override and score-formula fall-through) that would otherwise route unconditionally to the retrieval-augmented tier. The Tier 1 and Tier 2 outcomes are unchanged from \Cref{tab:dispatch-rules}: a query that earns the memory-direct or zero-shot dispatch on its own combined score does not consult the classifier.}
\label{tab:ruc-dispatch}
\end{table}

\begin{algorithm}[!htbp]
\caption{Retrieval-utility classifier scoring and dispatch refinement.}
\label{alg:ruc}
\begin{algorithmic}[1]
\Require question $q$, top-ranked passage $\mathbf{p}_{1}$, dense top-five passages $\mathbf{P}_{5}$, pre-routing decision $d_{\text{pre}} \in \{\text{Tier 1}, \text{Tier 2}, \text{Tier 3}\}$
\Require frozen scaler $\Sigma$, frozen logistic head $w$, threshold $\tau_{\text{RUC}} = 0.375$
\If{$d_{\text{pre}} \in \{\text{Tier 1}, \text{Tier 2}\}$}
    \State \Return $d_{\text{pre}}$ \Comment{classifier does not gate Tier 1 or Tier 2 outcomes}
\EndIf
\State $\mathbf{x}_{\text{q}} \gets$ question-surface features ($q$) \Comment{interrogative one-hot, negation, temporal, numerical, myth-regex, entity count}
\State $\mathbf{x}_{\text{p}} \gets$ top-one-passage features ($q, \mathbf{p}_{1}$) \Comment{similarity, entity overlap, answer-type matches}
\State $\mathbf{x}_{\text{r}} \gets$ retrieval-battery features ($\mathbf{P}_{5}$) \Comment{FAISS top-five statistics, cross-encoder rerank statistics, pairwise NLI statistics, Wikidata}
\State $p_{\text{IK}} \gets$ base-generator self-knowledge probe ($q$) \Comment{Kadavath-style probe under the disable-adapter context}
\State $\mathbf{x} \gets [\mathbf{x}_{\text{q}}, \mathbf{x}_{\text{p}}, \mathbf{x}_{\text{r}}, p_{\text{IK}}]$
\State $p_{\text{RAG}} \gets \sigma\!\left(w^{\top}\, \Sigma(\mathbf{x})\right)$ \Comment{logistic on standard-scaled features}
\If{$p_{\text{RAG}} \ge \tau_{\text{RUC}}$}
    \State \Return $\text{Tier 3}$ \Comment{retrieval-augmented generation}
\Else
    \State \Return $\text{Tier 2}$ \Comment{zero-shot generation}
\EndIf
\end{algorithmic}
\end{algorithm}

The contribution of the retrieval-utility classifier to the registered hypotheses is twofold. First, by routing retrieval-hurt queries to the zero-shot tier, the classifier closes a known regression mode on misconception-bait benchmarks that the architecture without the classifier exhibits and that the deferred-buffer and retroactive re-verification mechanisms cannot close because both consume the architecture's output rather than its tier choice. Second, by separating the retrieval-helpful subpopulation from the retrieval-hurt subpopulation at routing time, the classifier sharpens the self-improvement training pool's signal-to-noise ratio: the pool admits Tier 3 answers on queries where retrieval helps and Tier 2 answers on queries where retrieval would have hurt, so the fine-tune cycle absorbs grounded reasoning on the former without absorbing passage-amplified errors on the latter. The empirical receipt for both contributions, including the realised TruthfulQA recovery, the per-cycle routing-distribution trajectory, and the head-to-head ablation against the architecture without the classifier on the same six-cycle trajectory, is reported in \Cref{sec:ruc-ablation}.

\subsection{Verifier ensemble}
\label{sec:verifier}

The verifier ensemble is the post-generation block that converts a generated answer into a vector of nine complementary signals across four families. The four-family decomposition follows the white-box-versus-black-box and instance-versus-aggregate axes catalogued in the recent confidence-estimation survey for language models \cite{geng2024survey}, but extends them with the explicit grounding-versus-relevance separation that the off-topic-but-grounded failure mode requires. The ensemble is the single source of every quantity that the calibrated probability composite consumes, and its design is governed by two principles: each signal must measure a different failure mode from a different evidence source, and no signal may be redundant with another in the same family. The nine signals together feed the per-signal isotonic regression and the regularised logistic boost layer specified in \Cref{sec:composite}; they also feed the four-outcome decision tree's confabulation early-exit rule specified in \Cref{sec:decision}. \Cref{tab:verifier-signals} summarises every signal, its family, its evidence source, the resampling configuration that generates it, and the failure mode it closes. \Cref{alg:verifier} states the ensemble's execution order as implemented in the unified verifier.

\begin{table}[!htbp]
\centering
\footnotesize
\begin{tabular}{@{}p{2.4cm}p{2.6cm}p{3.0cm}p{2.6cm}p{3.4cm}@{}}
\toprule
\textbf{Signal} & \textbf{Family} & \textbf{Evidence source} & \textbf{Resampling} & \textbf{Failure mode closed} \\
\midrule
$u_{\text{token}}$        & Internal calibration & Generator log-probs                    & Single decode          & Token-level low-confidence generation \\
$u_{\text{dropout}}$      & Internal calibration & MC-dropout in train mode               & $K{=}5$ passes         & Hidden-state instability \\
$u_{\text{internal}}$     & Internal calibration & Composite of $u_{\text{token}}$, $u_{\text{dropout}}$ & Derived         & Internal-confidence summary \\
$s_{\text{avg}}$          & Sample-set agreement & Generator resampled chains             & $M{=}3$ at $T{=}0.7$   & Disagreement across chains \\
$p_{\text{entail}}$       & Sample-set agreement & NLI judge on chain $\to$ answer        & Averaged over $M$      & Chain-to-answer support \\
$p_{\text{ground,max}}$   & External grounding   & NLI judge on passage $\to$ answer      & Top-1 reranked         & No supporting passage \\
$p_{\text{ground,mean}}$  & External grounding   & NLI judge on passage $\to$ answer      & Mean over top-3        & Weak average grounding \\
$p_{\text{ground,atomic}}$ & External grounding  & NLI judge on atomic facts              & Weakest link, length $\ge 12$ tokens & Partial-support fabrication on long answers \\
$q_{\text{a-rel}}$        & QA relevance         & Cross-encoder on (question, answer)    & Single scoring         & Off-topic answer with matching passage \\
\bottomrule
\end{tabular}
\caption{The nine verifier signals that enter the calibrated probability composite, organised across four families. Resampling counts and length gate match the shipped verifier implementation; no fold is shared with the calibration fold.}
\label{tab:verifier-signals}
\end{table}

\begin{algorithm}[!htbp]
\caption{Verifier ensemble (Stage~5).}
\label{alg:verifier}
\begin{algorithmic}[1]
\Require query $q$, answer $a$, generator $G$, NLI judge $J$, cross-encoder $X$, retrieved passages $\mathcal{P}$
\Statex \textbf{Internal calibration family}
\State $u_{\text{token}} \gets$ mean token log-prob of $a$ under $G$ at the generation temperature
\State $u_{\text{dropout}} \gets$ variance of token log-prob of $a$ across $K{=}5$ MC-dropout passes through $G$
\State $u_{\text{internal}} \gets 0.5\,u_{\text{token}} + 0.5\,(1-u_{\text{dropout}})$
\Statex \textbf{Sample-set agreement family}
\State $\{a^{(m)}\}_{m=1}^{M=3} \gets$ resample chains from $G$ at temperature $0.7$
\State $s_{\text{avg}} \gets$ mean unique-pair cosine similarity across $\{a^{(m)}\}$ in the sentence-embedding space
\State $p_{\text{entail}} \gets \tfrac{1}{M}\sum_{m=1}^{M} J(a^{(m)} \to a)$
\Statex \textbf{External grounding family}
\State $\{e_{p}\}_{p\in \mathrm{top}{-}3(\mathcal{P})} \gets J(p \to a)$ for each reranked passage $p$
\State $p_{\text{ground,max}} \gets \max_{p} e_{p}$
\State $p_{\text{ground,mean}} \gets \mathrm{mean}_{p}\, e_{p}$
\If{$|a| \ge 12$ tokens}
    \State decompose $a$ into atomic facts $\{f_{i}\}$
    \State $p_{\text{ground,atomic}} \gets \min_{i}\, \max_{p}\, J(p \to f_{i})$
\Else
    \State $p_{\text{ground,atomic}} \gets p_{\text{ground,mean}}$ \Comment{length-gate fallback}
\EndIf
\Statex \textbf{Question-answer relevance}
\State $q_{\text{a-rel}} \gets X(q, a)$ \Comment{cross-encoder scoring}
\State \Return $(u_{\text{token}}, u_{\text{dropout}}, u_{\text{internal}}, s_{\text{avg}}, p_{\text{entail}}, p_{\text{ground,max}}, p_{\text{ground,mean}}, p_{\text{ground,atomic}}, q_{\text{a-rel}})$
\end{algorithmic}
\end{algorithm}

\subsubsection{Internal calibration signals}
\label{sec:verifier-internal}

The internal-calibration family reads the generator's own state on the answer it produced. The token-level uncertainty signal is the mean log-probability of the answer's tokens under the generation, with higher values reading as higher confidence; this signal captures the most local form of generator self-doubt. The dropout-disagreement signal is the variance of the answer's log-probability across the registered $K=5$ Monte-Carlo dropout passes through the generator with dropout layers active, with lower variance reading as higher confidence; this is the Bayesian-approximation interpretation of dropout introduced by \cite{pmlr-v48-gal16}, and it captures the stability of the generator's hidden state under stochastic ablation. The internal-state composite signal aggregates the two preceding signals into a single scalar with the fixed combination $0.5\,u_{\text{token}} + 0.5\,(1-u_{\text{dropout}})$, registered before the calibration phase; this signal acts as the input to the confabulation early-exit rule because it isolates the generator's internal-confidence channel from the external-grounding channel. The internal-calibration family is the only family that reads the generator's hidden state directly and is therefore the only family whose distributions shift after every cycle's fine-tune; this drift is what motivates the retroactive re-verification pass in \Cref{sec:retroverify}. The hidden-state-eigenvalue alternative \cite{chen2024inside}, which scores the log-determinant of the per-query response covariance, was considered as a candidate replacement for the dropout signal but was not adopted in the locked configuration for cost-efficiency reasons.

\subsubsection{Sample-set agreement signals}
\label{sec:verifier-sample-set}

The sample-set family probes the generator's stability by resampling answers under temperature and measuring agreement across the resampled chains. The semantic-consistency signal is the mean unique-pair cosine similarity across $M=3$ chains at temperature $T=0.7$ in the sentence-embedding space, with higher values reading as higher agreement \cite{wang2023selfconsistency}. The chain-to-answer entailment signal is the mean entailment of each of the $M$ resampled chains against the served answer under the natural-language-inference judge, with higher entailment reading as higher support; this signal closes the failure mode where the served answer is well-formed but is not the conclusion that the model's own resampling supports. The two signals share the same $M=3$ resampling pool, so semantic-consistency and chain-to-answer entailment are amortised against a single decode-phase cost. The family adds non-trivial decode-phase work to the per-query cost, which is why pre-routing routes the query to the memory-direct tier when it can, deferring the sample-set work to queries that genuinely need it.

Empirical retirements at the sample-set family. Two semantic-dispersion signals were registered earlier in the project and retired before the main run on the registered cycle-zero calibration evidence. The normalised semantic entropy signal of \cite{farquhar2024semantic,kuhn2023semantic}, computed over $K=10$ chains at temperature $T=1.0$, registered a Pearson correlation indistinguishable from zero with exact-match labels on every per-benchmark slice of the calibration fold and a boost-layer weight within five times ten to the negative four of zero on every fit. The retained signal contributed no discriminative information beyond what the semantic-consistency and chain-to-answer entailment signals on the same $M=3$ chains already provided, and the $K=10$ stochastic generations needed to compute it dominated the verifier's wall-clock cost. The kernel-language-entropy generalisation that uses graded semantic-similarity kernels rather than hard cluster assignments \cite{nikitin2024kle} was registered as a candidate upgrade for the signal but was not adopted in the shipped configuration. Both the entropy signal and the kernel-entropy upgrade are retired in the registered architecture; the deployed verifier consumes nine signals across the four families and the calibrated probability composite ingests the same nine.

\subsubsection{External grounding signals}
\label{sec:verifier-grounding}

The external-grounding family scores the answer against the retrieved passages under three aggregation rules. The top-passage entailment signal is the entailment score of the highest-scoring passage against the answer under the natural-language-inference judge, with higher entailment reading as stronger support; this signal captures the case where a single retrieved passage strongly supports the answer. The pooled-mean entailment signal is the mean entailment over the top-three reranked passages, with higher mean reading as broader support; this signal captures the case where multiple passages each contribute partial support and the answer is the integration over them. The atomic-fact entailment signal decomposes the answer into atomic facts and computes the weakest-link maximum entailment across the top reranked passages for each atomic fact, with the final score being the minimum across the atomic facts; this signal captures the case where an answer is mostly supported but contains a single fabricated detail that the pooled-mean signal would average away. The atomic-fact signal fires only on answers that exceed a length gate of twelve tokens, registered in the implementation, because the atomic decomposition of a short label answer reduces to the answer itself and adds noise rather than signal. Below the gate the atomic signal degrades to the pooled-mean signal, so that a shorter answer is not penalised by an incoherent decomposition.

\subsubsection{Question-answer relevance}
\label{sec:verifier-qa-rel}

The question-answer relevance signal scores the served answer against the original question under a cross-encoder, with higher score reading as higher topical relevance. This signal is treated as a separate signal rather than folded into the grounding family because the failure mode it closes is structurally different: an off-topic answer can be perfectly grounded in a passage that is itself off-topic for the question, and the grounding family alone would not detect this case. The cross-encoder is shared with the reranker on the retrieval-augmented tier, so the per-query cost of this signal is small relative to the work the cross-encoder is already doing on the same query. The contribution to the registered hypotheses is direct: the question-answer relevance signal is what closes the off-topic-but-grounded failure mode that the composite hallucination metric's off-topic subtype tracks, and it is one of two signals (alongside the chain-to-answer entailment) whose calibrated weight in the boost layer is expected to be non-trivial under the empirical fits reported in \Cref{ch:evaluation}.

\subsection{Calibrated probability composite}
\label{sec:composite}

The calibrated probability composite is the function that converts the nine verifier signals into a single calibrated storage score that the fixed-threshold gate consumes. The composite is a three-stage pipeline. The first stage is per-signal isotonic regression that maps each raw signal value to a calibrated probability of correctness, fitted independently per signal on a disjoint calibration fold. The second stage is a per-benchmark refinement that fits a child composite on each training-benchmark slice of the calibration fold and blends each child's monotone curve with the pooled fit through a James-Stein style shrinkage prior, so that benchmarks with weaker per-bench evidence regress toward the pooled fit rather than overfitting on small per-bench folds. The third stage is a regularised logistic boost layer that aggregates the per-signal calibrated probabilities through a weighted log-odds sum, with the weights and the intercept fitted on the same calibration fold against exact-match labels. \Cref{alg:composite} states the fit-time and predict-time procedures as implemented; the locked configuration parameters used in the main run are pulled from the cycle-zero composite calibration artefact.

\begin{algorithm}[!htbp]
\caption{Calibrated probability composite (fit and predict).}
\label{alg:composite}
\begin{algorithmic}[1]
\Statex \textbf{Fit stage.} Given a calibration fold partitioned by training benchmark $\{(\mathcal{D}_{b})\}_{b\in\mathcal{B}}$ with each $\mathcal{D}_{b}$ a set of pairs $(x_{i}, y_{i})$ where $x_{i}\in\mathbb{R}^{9}$ is the per-signal vector and $y_{i}\in\{0,1\}$ is the exact-match label; let $\mathcal{D}_{\text{pool}} = \bigcup_{b}\mathcal{D}_{b}$ be the pooled fold.
\Statex \emph{Pooled fit.}
\For{each signal $j \in \{1,\dots,9\}$}
    \State direction $\gets \mathrm{sgn}\!\left(\mathrm{Pearson}(x_{:,j}, y)\right)$ on $\mathcal{D}_{\text{pool}}$ \Comment{auto-detect monotone direction}
    \State $\mathrm{iso}^{\text{pool}}_{j} \gets \mathrm{IsotonicRegression}(\text{increasing}=\text{direction}, y_{\min}{=}0.02, y_{\max}{=}0.98)$
    \State $\mathrm{iso}^{\text{pool}}_{j}.\mathrm{fit}(x_{:,j}, y)$ on $\mathcal{D}_{\text{pool}}$
\EndFor
\Statex \emph{Per-benchmark fit with shrinkage prior.}
\For{each benchmark $b \in \mathcal{B}$ and each signal $j$}
    \State fit a child isotonic curve $\mathrm{iso}^{b}_{j}$ on $\mathcal{D}_{b}$ alone
    \State at every knot $x$ of $\mathrm{iso}^{b}_{j}$, blend the calibrated probability with the pooled curve, $\mathrm{iso}^{b}_{j}(x) \gets \alpha\cdot\mathrm{iso}^{b}_{j}(x) + (1-\alpha)\cdot\mathrm{iso}^{\text{pool}}_{j}(x)$, with the locked shrinkage strength $\alpha=0.6$
\EndFor
\Statex \emph{Boost layer.}
\State Build feature matrix $L \in \mathbb{R}^{N\times 9}$ with $L_{i,j} = \log\!\left(\tfrac{\mathrm{iso}^{\text{pool}}_{j}(x_{i,j})}{1-\mathrm{iso}^{\text{pool}}_{j}(x_{i,j})}\right)$ on the pooled fold
\State $(\mathbf{w}, b) \gets \mathrm{LogisticRegression}(C{=}0.01,\ \text{penalty}=\ell_{2},\ \text{max\_iter}{=}2000).\mathrm{fit}(L, y)$
\State Repeat the boost fit on each per-benchmark fold using the corresponding shrunk isotonic curves to produce per-benchmark weights $(\mathbf{w}_{b}, b_{b})$
\Statex \textbf{Predict stage.} Given a query's per-signal vector $x \in \mathbb{R}^{9}$ and benchmark tag $b^{\star}$:
\State if $b^{\star}\in\mathcal{B}$ select per-benchmark $(\mathrm{iso}^{b^{\star}}_{j},\, \mathbf{w}_{b^{\star}},\, b_{b^{\star}})$, else fall back to the pooled $(\mathrm{iso}^{\text{pool}}_{j},\,\mathbf{w},\,b)$
\State $\ell \gets b + \sum_{j=1}^{9} w_{j}\cdot\log\!\left(\tfrac{\mathrm{iso}_{j}(x_{j})}{1-\mathrm{iso}_{j}(x_{j})}\right)$
\State $u_{\text{stored}} \gets \sigma(\ell) = 1 / (1 + \exp(-\ell))$
\State \Return $u_{\text{stored}}$
\end{algorithmic}
\end{algorithm}

\subsubsection{Per-signal isotonic regression}
\label{sec:composite-isotonic}

Per-signal isotonic regression is preferred over the alternatives at this site for three reasons. Isotonic regression admits a non-linear monotone mapping rather than the single-scalar log-odds shift of Platt scaling, which matters because the verifier signals do not all admit a logistic relationship to correctness; the top-passage entailment signal in particular saturates at the upper end of its range, and Platt scaling would absorb that saturation into a globally too-soft probability assignment. Isotonic regression is data-efficient relative to histogram binning at the calibration-fold size used in this thesis, because it shares mass across adjacent bins through its monotone structure rather than partitioning the fold into independent bins that each need their own minimum sample count. Isotonic regression auto-detects the monotone direction from the Pearson correlation between the raw signal and the binary label, so a signal that reads high-confidence-as-low-value (such as the dropout-disagreement signal, where higher variance reads as lower confidence) is calibrated correctly without a registry of per-signal directions \cite{niculescu2005predicting}. The implementation clips the calibrated probability into the range $[0.02,\,0.98]$ to keep the log-odds finite and to bound the influence of any single signal on the composite. Each signal must contribute at least fifty calibration-fold samples for its isotonic fit to be admitted; below this floor the signal is dropped and its log-odds contribution defaults to zero in the composite, so that the composite degrades gracefully when a fold is too small.

\subsubsection{Per-benchmark fit with shrinkage prior}
\label{sec:composite-shrinkage}

The composite is fitted at two granularities on the same calibration fold: a pooled fit over the union of training-benchmark slices, and a per-benchmark fit over each slice separately. The pooled fit serves as both the fallback for any query whose benchmark tag is not in the registered training panel and as the prior toward which each per-benchmark child curve is regularised. The shrinkage rule is the James-Stein style convex blend, applied at every knot of the per-benchmark curve, that replaces the child knot value with a weighted combination of the child's own value and the pooled curve evaluated at the same input. The locked shrinkage strength is $\alpha=0.6$, which keeps a moderate per-benchmark adaptation while pulling each child curve away from the per-bench overfit that pure local fitting produces on weak-signal panels. The shrinkage strength was selected by the cycle-zero diagnostic that fitted the pure per-bench composite ($\alpha=1$) and observed that the per-benchmark area-under-the-receiver-operating-characteristic readings on the weak-signal training panels regressed below the pooled reading by ten to fifteen points, signalling that the per-bench fold size of approximately five hundred samples per benchmark was too small to support the more flexible per-bench isotonic envelope unaided. The blend at $\alpha=0.6$ moves the regression-AUROC readings back into parity with the pooled fit on those panels while preserving a measurable per-bench improvement on the strong-signal benchmarks. The per-benchmark child also retains its own boost coefficients fitted on the bench-local fold, so the per-bench curve and the per-bench weights jointly capture the benchmark's own reliability profile rather than borrowing the pooled boost weights uniformly.

\subsubsection{Boost layer with L2 regularisation}
\label{sec:composite-boost}

The boost layer aggregates the per-signal calibrated probabilities by fitting a logistic regression on the per-signal log-odds as features against the exact-match labels of the calibration fold \cite{cherian2024conformal}. The learned coefficients become the per-signal weights and the intercept becomes the composite's bias term; at predict time the weighted log-odds sum is sigmoid-back-projected to a probability in the unit interval. A learned linear combination over the per-signal log-odds is preferred over an identity-weighted sum at this site because the verifier signals carry overlapping information across families: the chain-to-answer entailment signal and the semantic-consistency signal both read sample-set agreement, the top-passage and pooled-mean entailment signals both read external grounding, and the boost layer is what allows the composite to absorb this redundancy by down-weighting the redundant component of each signal's contribution. The L2 penalty is preferred over L1 at this site because every signal in the registered set carries non-zero discriminative signal under the registered hypotheses; the L2 penalty shrinks weights smoothly without setting any to zero, which preserves the architectural commitment to the nine-signal coverage and keeps the boost layer a reweighting rather than a feature-selection step. The locked penalty strength is $C = 0.01$, which corresponds to a strong L2 regulariser; this value was registered in the cycle-zero calibration sweep against the alternatives $C\in\{0.01,\,0.10,\,1.00\}$, and the empirical evidence for the strong-regulariser choice is reported under the calibration sweep panel in \Cref{ch:evaluation}. The boost layer is fitted in parallel for the pooled and per-benchmark composites so that the per-benchmark child carries its own per-signal weights, capturing the case where the same raw signal reads in a different direction on a different benchmark; the pooled boost coefficients are used as the fallback whenever the verifier sees a query whose benchmark tag is not in the registered panel.

\subsection{Fixed-threshold storage gate}
\label{sec:conformal-gate}

The storage gate is the procedure that converts the calibrated storage score from \Cref{sec:composite} into a storage decision at a pre-registered cut-point on the calibrated probability scale. The gate is a fixed-threshold rule: an entry is admitted to the storage class whenever its calibrated composite satisfies the storage condition, deferred whenever it satisfies the deferral condition, and otherwise routed to abstention or discard by the rule specified in \Cref{sec:decision}. The fixed-threshold formulation replaces the split-conformal formulation considered earlier in the project, which fitted the storage and deferral thresholds at pre-registered precision targets and re-fit them at every cycle boundary under exchangeability between the calibration fold and the live distribution. The empirical evidence accumulated through the cycle-zero calibration sweep showed that the calibration-versus-evaluation distribution shift on the five-benchmark panel breaks the marginal-coverage premise of the split-conformal procedure, with the realised eval-fold precision falling fifteen to fifty percentage points below the conformal target on training benchmarks and the per-benchmark fits collapsing to the rejection limit on the weak-signal panels. The fixed-threshold replacement is preferred over the split-conformal alternative on two grounds: it is robust to the cal-versus-eval shift that the conformal premise requires, in the lineage of post-hoc calibration analyses by \cite{niculescu2005predicting} and the selective-classification framing of \cite{geifman2017selective}; and it commits the gate's operating point to a single registered scalar that any reader can locate in the configuration without reconstructing the conformal fit.

The gate's two cut-points are registered in the configuration before the main run begins. The storage threshold $\tau_{\text{store}} = 0.60$ is the cut-point at which the storage class is admitted; the deferral threshold $\tau_{\text{defer}} = 0.45$ is the cut-point below which an entry is held in the deferral buffer pending cycle-boundary re-evaluation; the abstention floor $\phi_{\text{pg}} = 0.20$ on the top-passage entailment signal separates the abstention outcome from the discard outcome below the deferral threshold. All three constants are read at every decision call from the configuration object, so a future re-tuning propagates without re-instantiating the verifier. The empirical calibration-fold precision at each cut-point is recorded by an empirical-precision audit that runs once per cycle on the cycle's labelled calibration fold and writes the audit table to the cycle's output directory; the audit is the empirical receipt for the calibration-consistent floor invoked by \Cref{thm:purity}.

\subsubsection{Storage threshold}
\label{sec:conformal-store}

The storage threshold is the operating point at which the storage class is admitted; an entry is admitted whenever its calibrated composite satisfies $u_{\text{stored}} \ge \tau_{\text{store}} = 0.60$. The choice of cut-point is registered against the cycle-zero threshold sweep, which compared two candidates on the labelled cycle-zero calibration fold: a stricter cut at $0.65$ admitted seventy correct entries out of one hundred and five total admissions for an empirical precision of approximately sixty-seven percent, while the registered cut at $0.60$ admitted three hundred and six correct entries out of four hundred and seventy-eight total admissions for an empirical precision of approximately sixty-four percent. The two operating points sit on the same precision shelf at the cycle-zero distribution, but the lower cut admits roughly four-and-a-half times as many correct memories at the same per-store precision and is the registered configuration for the main run. The storage class is what populates the episodic memory and what the training pool draws from at each cycle boundary, so the gate is paired with the strict training threshold above the storage threshold and with the retroactive prune threshold below it; the storage cut alone is not the only barrier between a candidate answer and the next cycle's gradient update.

\subsubsection{Deferral threshold}
\label{sec:conformal-defer}

The deferral threshold defines the catchment for the deferral buffer; an entry is buffered whenever its calibrated composite falls in the band $\tau_{\text{defer}} \le u_{\text{stored}} < \tau_{\text{store}}$. The registered cut-point is $\tau_{\text{defer}} = 0.45$. The looser cut-point is chosen because the deferral class does not propagate to the training pool at the next cycle boundary; deferred entries are revisited under the updated generator at the cycle-boundary deferred-entry reconsideration pass specified in \Cref{sec:deferred-reconsider} and either promoted into the storage class or expired by the time-to-live rule. Setting the deferral threshold loosely admits a wider band of marginal entries into the buffer, where they can be re-evaluated under stronger evidence at the next cycle boundary, while keeping the storage class itself disciplined by the stricter storage cut. The decision tree that consumes both cuts enforces the ordering invariant $\tau_{\text{defer}} \le \tau_{\text{store}}$ structurally; the configuration object refuses to load a state in which the two values would be swapped.

\subsubsection{Cycle-boundary recalibration}
\label{sec:conformal-cycle-recalib}

The locked configuration draws a sharp line between architectural parameters and calibration parameters. Architectural parameters stay frozen across every cycle of the main run: the storage and deferral cut-points $\tau_{\text{store}} = 0.60$ and $\tau_{\text{defer}} = 0.45$, the abstention floor $\phi_{\text{pg}} = 0.20$, the retroactive prune threshold $\tau_{\text{retro}} = 0.50$, and the strict training threshold above the storage cut. Calibration parameters re-fit at every cycle boundary under a single composite-recalibration pass that consumes the cycle's freshly-scored calibration fold and writes the per-signal isotonic curves, the per-benchmark child curves, the shrinkage-blended knot tables, and the boost layer's weights and intercept to the cycle's output directory. The recalibrated composite is then promoted to the canonical path that the verifier reads on the next pipeline initialisation, so the cycle-$N$ verifier scores entries through cycle-$N$ calibration rather than the previous cycle's stale curves.

The composite refit is sequenced inside the cycle boundary so that the post-fine-tune calibration governs the retroactive re-verification pass rather than the previous cycle's stale calibration. After the self-improvement loop commits its weight update, the calibration fold is rescored under the now-improved generator, the recalibration script re-fits the per-signal isotonic curves and the boost coefficients against this freshly-scored fold, and the verifier is rebound to the cycle-$N$ calibration block before retroactive re-verification reads through it. Without this ordering the retroactive pass would score memory entries through the previous cycle's isotonic envelope whose fit no longer covers the post-fine-tune raw signal distribution, and exact-match-correct stored entries with high post-fine-tune token confidence would map to artefactually-low calibrated probabilities and trip the prune threshold; the corrected sequence holds the retroactive pass under the cycle's own freshly-fitted curves and avoids that artefact. The empirical-precision audit at the locked cut-points runs once per cycle against the same recalibrated fold and writes the audit table that the calibration-consistent floor invoked by \Cref{thm:purity} reads.

The split between frozen architectural cut-points and adaptive composite calibration preserves the architectural anchor that Hypothesis~1 of \Cref{sec:hypotheses} tests: the nine-signal composite design is the object under empirical test across the trajectory, and re-tuning the cut-points cycle-by-cycle would conflate verifier-architecture quality with operating-point drift correction. The split also surfaces the research-versus-production parity: production deployment runs the same cycle-boundary composite refit as the research trajectory, with the cycle horizon set by accumulation rather than by a fixed query-budget schedule. A production cycle ends when a registered number of storage-class admissions has accumulated, or on a calendar trigger, whichever is sooner; at the cycle boundary the same labelled calibration fold is constructed (a stratified sample of the cycle's accumulated admissions, sent for offline labelling), the same per-signal isotonic curves and per-benchmark children and shrinkage blend and boost coefficients are refitted on that fold, and the same retroactive re-verification pass runs against the freshly refitted verifier. The production pipeline therefore inherits the cycle-boundary calibration discipline from research; the only operational difference is the cycle-trigger policy and the labelling cadence, both of which are specified in the production runbook accompanying this thesis.

\subsection{Four-outcome storage decision}
\label{sec:decision}

The four-outcome decision tree converts the calibrated storage score and a small set of guard signals into one of four mutually exclusive actions: store, defer, abstain, or discard. The four outcomes are preferred over a binary store-or-discard rule for two reasons. First, a binary rule conflates the genuinely-uncertain case with the confidently-wrong case: a query whose composite falls just below the storage threshold is not the same as a query whose composite is near zero, and treating both as discard discards information that the next cycle's updated generator could resolve. The deferral outcome captures the genuinely-uncertain band into a bounded buffer that is revisited at the cycle-boundary deferred-entry reconsideration pass. Second, a binary rule does not separate the failure mode where the system has nothing to say from the failure mode where the system would say something but cannot ground it; the abstention outcome captures the explicit-refusal case as a first-class user-facing action rather than a silent discard, in the spirit of the conformal-abstention framework that bounds hallucination rate at a user-specified target \cite{yadkori2024abstention}.

The decision tree is governed by three thresholds and one early-exit rule. The two cut-points $\tau_{\text{store}} = 0.60$ and $\tau_{\text{defer}} = 0.45$ are inherited from \Cref{sec:conformal-gate} and govern the store and defer outcomes. A third threshold $\phi_{\text{pg}} = 0.20$ on the top-passage entailment signal separates the abstention outcome (no supporting passage) from the discard outcome (supporting passage but uncertain composite). The confabulation early-exit rule fires when the internal-confidence composite is high but the top-passage entailment is low,
\begin{equation}
\big(u_{\text{internal}} \ge 0.70\big) \;\wedge\; \big(p_{\text{ground,max}} \le 0.20\big)
\;\Rightarrow\; \texttt{DISCARD with early\_exit\_triggered=True},
\label{eq:early-exit-rule}
\end{equation}
and is checked before the storage and deferral cut-points. The early-exit rule is the architectural defence against the most damaging failure mode in the verifier ensemble: a generator that is internally confident but externally ungrounded. Without this rule the calibrated composite alone might assign such an answer a marginal storage score and admit it under the deferral threshold, where it would survive into the next cycle's reconsideration pass; the early-exit rule routes it directly to discard with a flag that the cycle-level diagnostics can read. \Cref{alg:decision} states the decision tree as implemented in the unified verifier.

\begin{algorithm}[!htbp]
\caption{Four-outcome storage decision (Stage~6).}
\label{alg:decision}
\begin{algorithmic}[1]
\Require storage score $u_{\text{stored}}$, internal-confidence composite $u_{\text{internal}}$, top-passage entailment $p_{\text{ground,max}}$
\Require fixed cut-points $\tau_{\text{store}} = 0.60,\ \tau_{\text{defer}} = 0.45$ from \Cref{sec:conformal-gate}
\Require constants $\eta_{u} = 0.70$, $\eta_{g} = 0.20$, $\phi_{\text{pg}} = 0.20$
\If{$u_{\text{internal}} \ge \eta_{u}$ \textbf{and} $p_{\text{ground,max}} \le \eta_{g}$}
    \State \Return $(\texttt{DISCARD},\ \texttt{early\_exit\_triggered}=\text{True})$ \Comment{confabulation early-exit}
\EndIf
\If{$u_{\text{stored}} \ge \tau_{\text{store}}$}
    \State \Return $(\texttt{STORE},\ \texttt{early\_exit\_triggered}=\text{False})$
\EndIf
\If{$u_{\text{stored}} \ge \tau_{\text{defer}}$}
    \State \Return $(\texttt{DEFERRED},\ \texttt{early\_exit\_triggered}=\text{False})$
\EndIf
\If{$p_{\text{ground,max}} < \phi_{\text{pg}}$}
    \State \Return $(\texttt{ABSTAIN},\ \texttt{early\_exit\_triggered}=\text{False})$
\EndIf
\State \Return $(\texttt{DISCARD},\ \texttt{early\_exit\_triggered}=\text{False})$
\end{algorithmic}
\end{algorithm}

The retroactive-review motivation that justifies the deferred path is structural rather than rhetorical. A deferred entry is rescored at the next cycle boundary under the updated generator (which is fitted on the pool that the current cycle's storage outcomes contributed to), so a deferred entry whose composite falls just below the storage threshold this cycle has a non-trivial chance of crossing the threshold next cycle as the generator's competence on its question class improves. Without the deferred path, every such entry would be permanently discarded after one cycle of borderline competence, and the fine-tune pool's growth rate would be bounded above by the storage rate at the strictest threshold the system was ever evaluated against. With the deferred path, the fine-tune pool grows monotonically as deferred entries are promoted, and the pool's quality is preserved because every promoted entry has cleared the storage threshold under the updated model rather than the older one. The deferred buffer is bounded in size by a registered capacity and entries that exceed a registered time-to-live without promotion are expired; the bound prevents the buffer from growing unboundedly while the cycle-boundary reconsideration pass clears most entries within a few cycles of admission.

\subsection{Cycle-boundary maintenance}
\label{sec:cycle-boundary}

The three cycle-boundary passes are also the architectural mechanisms that close the two statelessness problems registered in \Cref{sec:arch-opportunity}. The first statelessness problem is that standard retrieval-augmented generation pays the full retrieval-and-generation cost in every session because the model never absorbs what retrieval just taught it; the self-improvement loop closes this by drawing its training pool from stored episodes whose answers were originally produced by the retrieval-augmented tier and were certified by the verifier, so the fine-tune folds verified retrieval-augmented reasoning into the generator's parametric capacity and the next cycle's zero-shot tier serves a measurable share of the previously retrieval-dependent queries without retrieval. The second statelessness problem is that even when the model already knows the answer, the standard protocol re-generates it from scratch on every recurrence and offers no consistency guarantee across sessions; the verified episodic store closes this through similarity-indexed cached answers served at embedding-lookup cost, with the retroactive re-verification pass at every cycle boundary refreshing the stored composite under the locked verifier ensemble and the cycle-boundary-refit composite against the updated generator's outputs, so that the cached answer's quality tracks the model's current capability rather than freezing at the moment of admission. The deferred-entry reconsideration pass is the bridge between the two mechanisms: it revisits entries that were close to the storage cut-point at admission time but did not clear it, so that the same query can be re-admitted under the updated generator without rerunning retrieval. The three passes share a logical role: each runs once at the end of every cycle and each consumes the updated generator that the cycle's fine-tune produces. The retroactive re-verification pass re-scores every stored episode under the locked verifier ensemble against the updated generator's outputs and the cycle-boundary-refit composite, then prunes or refreshes accordingly. The deferred-entry reconsideration pass revisits every entry in the deferral buffer and either promotes, drops, or keeps the entry. The self-improvement loop draws a training pool from the storage class, runs the parameter-efficient fine-tune through the bounded low-rank-adapter envelope on the frozen backbone (the registered $L_{2}$ anchor coefficient is zero on this path; the full-fine-tune fallback for which the anchor is non-zero is registered but inactive in the shipped configuration), and either commits or rolls back the resulting weights against the multi-modal retention guard. All three passes are sequenced so that the fine-tune happens after the deferred-reconsider pass and before the retroactive re-verification pass for the next cycle, so that the retroactive pass always sees the freshly fine-tuned model rather than the model that produced the entries it is re-scoring.

\subsubsection{Retroactive re-verification}
\label{sec:retroverify}

The retroactive re-verification pass re-scores every stored episode under the current cycle's generator and applies four operations: prune entries whose reasoning chain trips the repetitive-loop filter, prune entries whose new storage score falls below the registered prune threshold, refresh the per-signal block of entries whose new score crosses the threshold, and update the stored composite to the new value. The first two operations remove entries that have lost their right to belong in memory; the third and fourth operations propagate the verifier's improved view of the entry forward to all downstream consumers, including the next cycle's training pool selection. The prune threshold is registered at $\tau_{\text{retro}} = 0.50$, below the deferral threshold so that an entry that is borderline-deferred-quality but no longer storage-quality is still kept for the deferred-reconsider pass; entries below the prune threshold are categorically removed. The full per-signal block (all nine verifier signals) is overwritten on every re-verification, not just the composite, so that the next cycle's calibration runs against the freshest verifier view of every stored episode. \Cref{alg:retroverify} states the pass.

\begin{algorithm}[!htbp]
\caption{Retroactive re-verification pass (Stage~8a).}
\label{alg:retroverify}
\begin{algorithmic}[1]
\Require memory $\mathcal{M}$ of stored entries, current verifier function $\mathrm{verify}(\cdot)$ under updated generator
\Require prune threshold $\tau_{\text{retro}} = 0.50$
\State $n_{\text{up}} \gets 0$, $n_{\text{rm}} \gets 0$
\For{each entry $e \in \mathcal{M}$}
    \If{$e.\mathrm{reasoning\_chain}$ trips the repetitive-loop filter}
        \State remove $e$ from $\mathcal{M}$ \textbf{; continue}
    \EndIf
    \State $\mathbf{x}_{\text{new}} \gets \mathrm{verify}(e)$ \Comment{full nine-signal re-scoring under updated model}
    \If{$\mathbf{x}_{\text{new}}.u_{\text{stored}} < \tau_{\text{retro}}$}
        \State remove $e$ from $\mathcal{M}$; $n_{\text{rm}} \gets n_{\text{rm}} + 1$
    \Else
        \State overwrite $e$'s per-signal block and stored composite with $\mathbf{x}_{\text{new}}$
        \State $n_{\text{up}} \gets n_{\text{up}} + 1$
    \EndIf
\EndFor
\State \Return $(n_{\text{up}},\, n_{\text{rm}})$
\end{algorithmic}
\end{algorithm}

\subsubsection{Deferred-entry reconsideration}
\label{sec:deferred-reconsider}

The deferred-entry reconsideration pass revisits every entry in the deferral buffer at every cycle boundary. Each entry is re-scored under the updated generator and the four-outcome decision tree is applied with the locked thresholds: an entry whose new score crosses the storage threshold is promoted into memory, an entry whose new score is below the deferral threshold is dropped, and an entry that remains in the deferral band is kept for the next cycle's reconsideration. A registered time-to-live of $\tau_{\text{ttl}} = 2$ cycles bounds the longest possible residency: an entry that survives two reconsideration passes without crossing the storage threshold is dropped categorically, so the buffer cannot accumulate entries that the model will never resolve. The buffer's hard capacity is registered at $K_{\text{def}} = 10\,000$ entries; when the cap is reached, the oldest entry is evicted to make room for new admissions. \Cref{alg:deferred-reconsider} states the pass.

\begin{algorithm}[!htbp]
\caption{Deferred-entry reconsideration pass (Stage~8b).}
\label{alg:deferred-reconsider}
\begin{algorithmic}[1]
\Require deferral buffer $\mathcal{B}$, current verifier $\mathrm{verify}(\cdot)$
\Require storage threshold $\tau_{\text{store}}$, deferral threshold $\tau_{\text{defer}}$, time-to-live $\tau_{\text{ttl}} = 2$
\State $n_{\text{pr}} \gets 0$, $n_{\text{dr}} \gets 0$, $n_{\text{kp}} \gets 0$
\For{each entry $e \in \mathcal{B}$}
    \State $\mathbf{x} \gets \mathrm{verify}(e)$
    \If{$\mathbf{x}.u_{\text{stored}} \ge \tau_{\text{store}}$}
        \State promote $e$ into the memory store; $n_{\text{pr}} \gets n_{\text{pr}} + 1$
    \ElsIf{$\mathbf{x}.u_{\text{stored}} < \tau_{\text{defer}}$ \textbf{or} $e.\mathrm{age}+1 \ge \tau_{\text{ttl}}$}
        \State drop $e$ from $\mathcal{B}$; $n_{\text{dr}} \gets n_{\text{dr}} + 1$
    \Else
        \State $e.\mathrm{age} \gets e.\mathrm{age} + 1$; keep $e$ in $\mathcal{B}$; $n_{\text{kp}} \gets n_{\text{kp}} + 1$
    \EndIf
\EndFor
\State \Return $(n_{\text{pr}},\, n_{\text{dr}},\, n_{\text{kp}})$
\end{algorithmic}
\end{algorithm}

The memory consolidation pass runs alongside the deferred-reconsider pass at every cycle boundary. It collapses near-duplicate stored episodes by clustering entries whose pairwise cosine similarity in the sentence-embedding space exceeds the consolidation threshold $\tau_{\text{cons}} = 0.92$, retaining the highest-quality representative within each cluster and merging the source-benchmark provenance of the dropped entries into the survivor. The threshold is set higher than the novelty threshold at write time, so that entries that survived the looser write-time filter but turn out to be near-paraphrases at consolidation time are still resolved into a single representative.

\subsubsection{Self-improvement loop}
\label{sec:sil}

The self-improvement loop is the cycle-boundary block that consumes the cycle's verified pool and produces the updated generator. It is an instance of the broader family of model-self-training procedures that fine-tune a generator on its own filtered outputs \cite{NEURIPS2022_639a9a17,hosseini2024vstar,gulcehre2023rest,ho2023large}, specialised here through the fixed-threshold storage gate that supplies the filter, the parameter-efficient adapter layer that absorbs the gradient update, the loss-reweighted training pool that prevents single-benchmark dominance, and the multi-modal retention guard that gates the commit. Its architectural role at the system level is the cost-amortisation mechanism that converts verified retrieval-augmented reasoning into zero-shot parametric capability. At cycle zero the storage class is dominated by entries that the retrieval-augmented tier produced, because the generator alone cannot yet answer most queries above the registered storage threshold. The training pool that the loop draws from these entries therefore consists of question-and-reasoning pairs whose correctness was established through retrieval. The fine-tune folds those pairs into the adapter, and the resulting generator at the next cycle is able to answer a measurable share of the same question distribution from the zero-shot tier without invoking retrieval. The redistribution of queries across the three serving tiers across cycles, reported empirically in \Cref{sec:tier-efficiency}, is therefore the operational reading of the loop's contribution: the increase in the zero-shot tier's share at the expense of the retrieval-augmented tier's share is the per-cycle measure of how much of the retrieval-mediated knowledge the generator has now internalised, and the cost saving on per-query latency is the deployment-side consequence. Recent theoretical work on whether such loops collapse or improve their generator establishes a sufficient condition that the verifier outperform the generator at the noisy-label-recovery task \cite{song2024gap}, and a complementary sufficient condition that the verifier-filtered admission policy be precision-bounded rather than recall-bounded so that the loop accumulates rather than displaces high-quality data \cite{fu2025collapse}; both conditions are met by the storage class admitted under the registered cut-points of \Cref{sec:conformal-gate}.

The loop has five operations, in order. First, sample selection from memory: the loop draws training pairs only from stored episodes whose composite score exceeds a strict training threshold registered above the storage cut-point, so the training pool is held to a higher bar than the user-facing retrieval tier. The selection is filtered through an in-distribution gate that excludes episodes whose source benchmark is disjoint from the registered training panel, so transfer-only benchmarks never feed back into the gradient update. Second, pool reweighting: the selected pairs are passed through a benchmark-balancing layer that prevents any single benchmark from dominating the gradient. The layer composes a temperature-mixed softmax over per-benchmark counts, a registered floor on the smallest per-benchmark contribution, a bounded upsampling cap on each benchmark's effective count, a hard ceiling on any single benchmark's share of the pool, and a cold-start fallback that injects gold-labelled pairs from the seed pool when a benchmark contributes too few verified episodes to clear the floor. The composition is iterative and produces a final pool whose share-by-benchmark distribution lies within the registered tolerance band of the target distribution. Third, the optimiser and the parameter envelope: the loop wraps the open-weight generator with a low-rank adaptation layer fitted on the attention and feed-forward projections, with a registered rank, scaling factor, and target-module list, and runs a thirty-two-bit decoupled-weight-decay AdamW under gradient checkpointing only on the adapter parameters at a learning rate ten times the backbone-fine-tune rate registered for the underlying generator. The full backbone weights are frozen across the cycle, which keeps the trainable footprint at roughly two percent of the backbone parameter count and bounds the cycle-to-cycle drift in the underlying generator's representation. The eight-bit optimiser path that would otherwise apply to a frozen-backbone-plus-full-fine-tune fallback is registered but inactive in the shipped configuration, because on the small adapter the per-parameter quantisation noise from eight-bit moment storage is relatively larger than the modest optimiser-state-memory saving against the thirty-two-gigabyte envelope. Fourth, the regulariser: the loss on the active low-rank-adapter path is cross-entropy on the verified training pool only,
\begin{equation}
\mathcal{L}(\boldsymbol{\theta}) \;=\; \mathcal{L}_{\mathrm{CE}}(\boldsymbol{\theta}) \;+\; \tfrac{\lambda_{\text{anc}}}{2}\, \big\| \boldsymbol{\theta} - \boldsymbol{\theta}^{(t-1)} \big\|_{2}^{2}\quad\text{with $\lambda_{\text{anc}}=0$ on the shipped adapter path,}
\label{eq:sil-anchor}
\end{equation}
where $\boldsymbol{\theta}$ is the adapter parameter vector, $\mathcal{L}_{\mathrm{CE}}$ is the cross-entropy loss on the verified training pool, and $\boldsymbol{\theta}^{(t-1)}$ is the previous cycle's adapter. The quadratic anchor term is registered as the parameter envelope for the inactive full-fine-tune fallback (when it engages, $\lambda_{\text{anc}}>0$ and the bracketed term becomes the uniform-Fisher approximation to elastic weight consolidation \cite{doi:10.1073/pnas.1611835114}). On the shipped low-rank-adapter path the bracketed term vanishes by design: the bounded parameter envelope of the adapter, layered on a frozen backbone, is itself the cycle-to-cycle drift bound, and an explicit weight-anchor would otherwise penalise deviation from a zero-anchored full-backbone reference whose weights the loop is not updating, which is meaningless arithmetic. Fifth, the retention guard: the loop probes a registered set of held-out distributions before the first cycle (the pristine anchor) and after every cycle, and computes the retention ratio (defined in \Cref{eq:nfr-retention-floor}) on each probe. The cycle commits the new adapter only when the worst probe ratio clears the floor $\rho_{\min} = 0.93$; otherwise the adapter is rolled back to the previous cycle's snapshot and the memory store is preserved unchanged. The probe set is multi-modal and intentionally diverse: a general-knowledge probe captures distributional drift away from the pretraining mass, while held-out test slices of the training-panel benchmarks capture in-task forgetting under the SIL gradient. The worst-probe rule is conservative on purpose: a loop that improves the average probe at the cost of regressing one specialised probe is the failure mode that Phase 1b's single-probe rule failed to catch. \Cref{alg:sil} states the loop.

\begin{algorithm}[!htbp]
\caption{Self-improvement loop with multi-modal retention guard (Stage~8c).}
\label{alg:sil}
\begin{algorithmic}[1]
\Require stored memory $\mathcal{M}$, previous adapter parameters $\boldsymbol{\theta}^{(t-1)}$, registered probe set $\mathcal{P}$
\Require training threshold $\tau_{\text{train}} > \tau_{\text{store}}$, anchor strength $\lambda_{\text{anc}}$, retention floor $\rho_{\min} = 0.93$
\State $\mathcal{T}_{0} \gets \{e \in \mathcal{M} : e.u_{\text{stored}} \ge \tau_{\text{train}}\}$ filtered by the in-distribution gate
\State $\mathcal{T} \gets$ pool-reweight$(\mathcal{T}_{0})$ \Comment{temperature softmax + floor + upsample + share cap + cold-start}
\State for each probe $p\in\mathcal{P}$: $\rho_{\text{pre}}^{p} \gets$ accuracy of $\boldsymbol{\theta}^{(t-1)}$ on $p$ (use cached pristine anchor when available)
\State $\boldsymbol{\theta}^{\star} \gets \arg\min_{\boldsymbol{\theta}}\, \mathcal{L}(\boldsymbol{\theta})$ from \Cref{eq:sil-anchor}, initialised at $\boldsymbol{\theta}^{(t-1)}$
\State for each probe $p\in\mathcal{P}$: $\rho_{\text{post}}^{p} \gets$ accuracy of $\boldsymbol{\theta}^{\star}$ on $p$
\If{$\min_{p\in\mathcal{P}}\, \rho_{\text{post}}^{p} / \rho_{\text{pre}}^{p} \ge \rho_{\min}$}
    \State $\boldsymbol{\theta}^{(t)} \gets \boldsymbol{\theta}^{\star}$ \Comment{commit the updated adapter}
\Else
    \State $\boldsymbol{\theta}^{(t)} \gets \boldsymbol{\theta}^{(t-1)}$ \Comment{rollback; memory preserved unchanged}
\EndIf
\State \Return $\boldsymbol{\theta}^{(t)}$
\end{algorithmic}
\end{algorithm}

The asymmetry between adapter and memory is the design property that makes the loop safe to iterate. A failed cycle costs the cycle's fine-tune energy but does not lose the verified episodes that the cycle accumulated; the next cycle inherits the previous adapter together with the larger memory, and the retroactive re-verification pass at the next cycle boundary will refresh the stored composites under the rolled-back adapter. The contribution to the registered hypotheses is direct: this asymmetry is what licenses the memory-growth invariant of \Cref{cor:stochastic-equilibrium}, because memory accumulates strictly across cycles even when the adapter does not.

\section{Theoretical analysis}
\label{sec:theory}

The architecture admits two complementary layers of theoretical guarantee that together characterise the trajectory of the storage class across cycles. The first layer treats the fixed-threshold storage gate as a calibration-consistent bound under the exchangeability assumption between the calibration fold and the cycle-$t$ population: the empirical precision of the calibration fold at the locked threshold projects forward as a one-sided lower bound on the cycle-$t$ pool's precision, supported by the post-hoc calibration analysis of \cite{niculescu2005predicting} and the selective-classification framing of \cite{geifman2017selective}. The second layer treats the storage gate as a Bayesian filter whose expected pool purity is determined by the verifier's balanced accuracy and the base generator's correctness rate, in the lineage of the recent analyses showing that retraining on verifier-correct labels strictly increases population accuracy in the binary linearly-separable setting \cite{das2025retraining}, that weak-to-strong pseudolabel training admits expansion-based generalisation bounds \cite{lang2024w2s}, and that the sharpening mechanism formalises the amortisation-via-self-training bound for language-model self-improvement \cite{huang2025sharpening}. To keep the receipts compact: a \emph{calibration-consistent} qualifier on a receipt indicates a finite-sample lower bound that holds at every cycle within the realised horizon under exchangeability between the cycle's calibration fold and its candidate pool; an \emph{asymptotic-projection} qualifier indicates an expected-value prediction whose trajectory shape is testable but whose limit is not directly observable within a short horizon.

\paragraph{Empirical limit of the realised trajectory:} The receipts below are evaluated against the six-cycle trajectory that closed at cycle five, when the equilibrium-saturation gate of \Cref{sec:env-cycles} terminated the run after the retention guard fired at cycle four and the post-rollback subsequence stabilised. This six-cycle window is the \emph{empirical limit} of the present thesis, not an asymptote: the theorems whose conclusions are infinite-horizon statements are evaluated against the cycle-five reading rather than against a fitted asymptote, and any residual gap between the realised limit and the theoretical asymptote is attributed to three implementation regularisers registered earlier in the chapter rather than to a falsification of the underlying mathematics. These regularisers are the $L_2$ adapter anchor of \Cref{eq:sil-anchor}, which pins the post-fine-tune weights near the previous-cycle checkpoint and prevents unbounded parametric drift; the multi-modal retention guard with rollback registered in \Cref{sec:retention}, which converts any cycle whose worst-probe ratio falls below the floor into a no-op on the parametric backbone; and the parametric capacity ceiling of the base generator and the dense passage index, which bound the pool's effective precision strictly below one even under an idealised verifier. The classification below preserves this distinction explicitly: theorems with cycle-local or bounded-band content are evaluated against the realised trajectory and reported with empirical receipts, theorems whose content is genuinely asymptotic are reframed against the empirical limit with the longer-horizon test deferred to \Cref{ch:conclusion}, and mathematical scaffolding identities that do not admit a thesis-side empirical contribution are gathered into a separate subsection at the end of the section so the reader can locate the formal machinery without confusing it with a falsifiable receipt. Proofs longer than half a page are sketched in the body and relocated to \Cref{app:reproducibility} in full.

\subsection{Pool purity}
\label{sec:theorem-purity}

Pool purity is the share of stored episodes whose stored answer is correct against the gold. \Cref{thm:purity} gives the calibration-consistent lower bound under the exchangeability precondition; the Bayes-identity expected value of \Cref{thm:bayes-purity} is recorded later in the section as design-rationale scaffolding because the realised receipt set evaluates the lower bound directly rather than the equality residual.

\begin{theorem}[Pool purity invariant: calibration-consistent lower bound on the pooled training-panel calibration fold]
\label{thm:purity}
Let $\mathcal{M}_{t}$ denote the storage pool at the end of cycle $t$, and let
\(\pi_{t} := \Pr\big(\mathrm{EM}=1 \,\big|\, e \in \mathcal{M}_{t}\big)\)
denote its precision under exact-match labelling, evaluated on the pooled training-panel calibration fold. Assume (i) the storage gate is the fixed-threshold rule that admits an entry whenever its calibrated composite satisfies $u_{\text{stored}} \ge \tau_{\text{store}}$, with $\tau_{\text{store}}$ frozen across cycles and $\pi_{\text{store}}$ the empirical calibration-fold precision at $\tau_{\text{store}}$ on a pooled training-panel fold exchangeable with the cycle-$t+1$ candidate pool at the pooled axis, (ii) the retroactive re-verification pass enforces the prune-threshold rule with $\tau_{\text{retro}} < \tau_{\text{store}}$ (any entry whose re-scored composite falls below $\tau_{\text{retro}}$ is removed; entries at or above $\tau_{\text{retro}}$ are retained, with the committed composite either raised or lowered to the new value), and (iii) the calibration-fold precision among entries that survive the retroactive pass is $\pi_{\text{retro}}$. Then
\begin{equation}
\pi_{t+1} \;\ge\; \min\!\big(\pi_{\text{store}},\,\pi_{\text{retro}}\big),
\label{eq:thm-purity}
\end{equation}
where $\pi_{\text{store}}$ and $\pi_{\text{retro}}$ are the calibration-fold expected precisions of the storage gate and the retroactive prune respectively, both reported on the pooled training-panel axis. The bound is stated on the pooled training-panel cal-fold axis only; per-benchmark slices of the candidate pool may violate the cal-fold-to-pool exchangeability premise on the per-bench axis, and the per-benchmark eval-fold $\pi_{\text{store}}$ is therefore reported as a within-scope-permitted diagnostic in the empirical receipt below rather than as a separate falsification target of \Cref{eq:thm-purity}.
\end{theorem}

\begin{proof}[Proof sketch]
At cycle $t+1$, $\mathcal{M}_{t+1}$ is the disjoint union of two sets: $\mathcal{M}_{t}^{\text{kept}}$, the entries from $\mathcal{M}_{t}$ that survived retroactive re-verification, and $\mathcal{N}_{t+1}$, the entries newly admitted by the storage gate. By assumption (ii) and (iii), every entry in $\mathcal{M}_{t}^{\text{kept}}$ has a re-scored composite at least $\tau_{\text{retro}}$, so its precision is bounded below by $\pi_{\text{retro}}$. By assumption (i) and exchangeability between the calibration fold and the candidate pool, the empirical precision of the calibration fold at $\tau_{\text{store}}$ projects forward as a lower bound on the precision of $\mathcal{N}_{t+1}$, which is therefore at least $\pi_{\text{store}}$. The precision of a union of disjoint sets is a convex combination of their precisions, bounded below by the smaller, giving \Cref{eq:thm-purity}.
\end{proof}

\paragraph{Empirical receipt (calibration-consistent floor):} Per-cycle pool purity reported in \Cref{ch:evaluation} should not drop below $\min(\pi_{\text{store}}, \pi_{\text{retro}})$ on the pooled training-panel calibration-fold axis the theorem registers. The pre-registered cycle-zero calibration-fold anchor at the locked threshold $\tau_{\text{store}}=0.60$ is $\pi_{\text{store}} \approx 0.64$, with the retroactive prune at $\tau_{\text{retro}}=0.50$ supplying the second floor; an empirical pooled-cal-fold pool-purity reading below the smaller of the two would falsify the exchangeability precondition or the retroactive-prune contract. Across the realised six-cycle trajectory the pooled cal-fold pool-purity reading remains in the range $0.72$ to $0.75$, clearing the registered floor by eight to eleven percentage points at every cycle, with the per-benchmark minimum on the pooled cal-fold axis of the training panel resting at $0.66$ on the TriviaQA stream chunk at cycle five and the retroactive-prune pass firing at every non-aborted cycle with a pruned share that declines monotonically across the committed-cycle subsequence from thirteen percent down to four percent. The bound is one-sided and re-evaluable at every cycle from the cycle's own pooled training-panel calibration fold. The receipt is a pool-level guarantee on the storage class on the pooled cal-fold axis and does not bound the held-out generation accuracy of the fine-tuned adapter; the cycle-three regression discussed in \Cref{sec:adj-cal-eval-gap} operates on the training-side gradient channel and is therefore outside the theorem's scope rather than a falsification of it.

\paragraph{Scope of the floor (per-benchmark eval-fold disclosure):} The floor that \Cref{eq:thm-purity} registers applies to the pooled training-panel calibration fold and not to per-benchmark eval-fold slices, because the cal-fold-to-pool exchangeability premise on which the projection rests is a pooled-axis premise that the per-bench axis does not inherit. The per-benchmark eval-fold reading reported in \Cref{tab:pi-store-trajectory} is therefore a within-scope-permitted diagnostic rather than a separate falsification target of the bound. On the realised trajectory the FEVER eval-fold $\pi_{\text{store}}$ dips below the registered $0.64$ floor at cycles one through four (readings in the band $[0.614, 0.634]$ at cycles one, two, three, and four respectively) and recovers to $0.649$ at cycle five, while the pooled cal-fold reading stays in the band $[0.732, 0.746]$ across the same committed cycles. The per-benchmark dip is consistent with FEVER's tighter precision regime under the locked configuration documented in \Cref{sec:adj-cal-eval-gap}: the per-benchmark composite refit under the shrinkage prior toward the pooled fit allocates per-benchmark slack toward the open-text factoid and five-option multiple-choice surfaces where the pre-routing confidence distribution is wider, and the FEVER per-benchmark eval-fold reading inherits the tighter operating point that the pooled axis absorbs. The reading is reported transparently because it is a real per-benchmark feature of the realised configuration, but it is not in scope for the floor of \Cref{eq:thm-purity} as stated above; the per-benchmark axis is registered as a future-work item in \Cref{ch:conclusion} alongside the longer-horizon receipts.

\subsection{Storage-rate sensitivity}
\label{sec:theorem-monotonicity}

The storage-rate sensitivity theorem characterises how the per-cycle admission count moves with the verifier's discrimination on the calibration fold. The companion identity on pool purity as a function of base accuracy is a textbook derivative of the Bayes identity and is recorded in the scaffolding subsection below.

\begin{theorem}[Monotonicity under generator improvement (verifier ensemble held fixed across cycles)]
\label{thm:monotone}
Let $d_{t}$ denote the calibration-fold class-separation Cohen's $d$ between em-correct and em-incorrect storage scores at cycle $t$, observed as a joint property of the cycle-$t$ generator's outputs on the cal fold, the locked verifier ensemble, and the cycle-$t$ refit composite; let $\sigma_{t}$ denote the empirical storage rate at the locked threshold $\tau_{\text{store}}$. The verifier ensemble (MiniCheck plus the long-hypothesis judge plus the signal-extraction layer) is held fixed across the trajectory; the only cycle-indexed components are the generator (advanced by the self-improvement loop on committed cycles) and the per-cycle composite refit (per-signal isotonic, per-benchmark shrinkage child fits, boost-layer weights, and intercept). Assume the calibrated storage score is approximately Gaussian within each correctness class with shared variance $\varsigma^{2}$ and class means $\mu_{+}(d_{t})$, $\mu_{-}(d_{t})$ separated by $d_{t}$. Let $p_{+}, p_{-}$ be the calibration-fold base rates of correct and incorrect outputs, and let $\phi_{\pm} = \phi\!\big((\mu_{\pm} - \tau_{\text{store}})/\varsigma\big)$ be the operating-point class densities. Then the sign of $\partial \sigma / \partial d$ equals the sign of $p_{+}\phi_{+} - p_{-}\phi_{-}$. In particular, when the operating-point correct-class mass dominates the incorrect-class mass and $d_{t+1} > d_{t}$, the storage rate is non-decreasing: $\sigma_{t+1} \ge \sigma_{t}$.
\end{theorem}

\begin{proof}[Proof sketch]
Under the within-class Gaussian assumption,
\(\sigma = p_{+} \Phi\!\big((\mu_{+} - \tau_{\text{store}})/\varsigma\big) + p_{-} \Phi\!\big((\mu_{-} - \tau_{\text{store}})/\varsigma\big)\),
where $\Phi$ is the standard normal CDF. Differentiating and using $\mu_{+} - \mu_{-} = d \cdot \varsigma$ yields
\(\partial \sigma / \partial d = (1/2)\big(p_{+} \phi_{+} - p_{-} \phi_{-}\big)\),
so the sign of $\partial \sigma / \partial d$ is the sign of $p_{+}\phi_{+} - p_{-}\phi_{-}$. The dominance condition holds whenever the threshold sits in the high-precision regime where the correct-class density exceeds the incorrect-class density at the operating point.
\end{proof}

\paragraph{Empirical receipt (storage rate in $d$, partial):} The deciding statistic is the sign of the per-cycle storage-rate change paired with the sign of the per-cycle Cohen's $d$ change on the calibration fold. Cycles for which the operating-point dominance condition $p_{+}\phi_{+} > p_{-}\phi_{-}$ holds are flagged in the trajectory artefact; a cycle in which the dominance condition does not hold is reported with its sign rather than asserted as a monotonicity violation. The realised cal-fold trajectory yields four observation points (cycles one, two, three, and five; cycle four's cal-fold pass was skipped by the retention-guard abort that rolled the adapter back to the cycle-three weights, see \Cref{sec:retention}). Of the three available transitions reported in \Cref{ch:evaluation}, the cycle-one to cycle-two transition sign-aligns with the theorem at large amplitude, the cycle-two to cycle-three transition shows a sub-noise Cohen's $d$ delta whose sign carries no information about the derivative, and the cycle-three to cycle-five transition shows a small reverse-sign storage-rate movement under increasing Cohen's $d$ with the dominance condition still holding at all four cycles. The reverse-sign transition is consistent with two confounds that the theorem's Gaussian model does not represent. The first is the per-cycle composite refit: the theorem treats the storage score as a fixed function whose only cycle-dependent quantity is the cal-fold class-mean separation $d_{t}$ (itself a downstream observable of the SIL-improved generator viewed through the locked verifier ensemble), whereas the implementation also refits the calibrated probability composite on each cycle's calibration fold, so the realised cycle-to-cycle movement in $\sigma_{t}$ bundles the generator-driven cal-fold discrimination shift with a composite-rescaling effect. The second is the shared-variance assumption: the empirical pooled-fold within-class standard deviation contracts across the trajectory rather than holding constant, so the proof's $\partial \sigma / \partial d$ decomposition is approximate cycle-to-cycle. The reported trajectory therefore evidences the sign-of-derivative prediction in the only large-amplitude transition observable on the cycle horizon and identifies the composite-refit and shared-variance confounds as the residual gap; a clean isolation would require freezing the composite weights and stratifying by class-variance estimate for an ablation cycle in which only the generator advances, which is registered as future work in \Cref{ch:conclusion}.

\subsection{Convergence}
\label{sec:theorem-convergence}

The convergence theorem characterises the per-cycle contraction of the precision gap to its long-run value. The realised six-cycle horizon is shorter than the horizon required to estimate a geometric rate cleanly, so the receipt is reframed as a bounded-band stationarity test on the committed-cycle subsequence and the sharper asymptotic rate is forwarded to future work.

\begin{theorem}[Convergence envelope: bounded-band stationarity on the committed-cycle subsequence]
\label{thm:convergence}
Let $\pi_{\infty} := \lim_{t \to \infty} \pi_{t}$ denote the asymptotic precision of the storage class, and let $G_{t} := |\pi_{t} - \pi_{\infty}|$ denote the cycle-$t$ gap. Under the assumptions of \Cref{thm:purity} and \Cref{thm:monotone} and the additional stationarity assumption that the per-cycle expected admission count $\sigma$ and the cal-fold class-separation $d$ (the joint observable of the cycle-$t$ generator, the locked verifier ensemble, and the cycle-$t$ refit composite) stabilise over the trajectory, there exists a contraction rate $r \in (0,1)$ such that, in expectation on the committed-cycle subsequence,
\begin{equation}
\mathbb{E}[G_{t}] \;\le\; (1-r)^{t} \, G_{0},
\label{eq:thm-convergence}
\end{equation}
where $\sigma$ is the per-cycle storage rate and $d$ is the verifier's discrimination on the calibration fold. The thesis-side receipt registered against \Cref{eq:thm-convergence} is the bounded-band stationarity reading on the realised committed-cycle subsequence, namely that $\pi_t$ remains within a band whose width is small relative to the per-cycle binomial sampling noise; the existence-claim form of $r \in (0,1)$ is what the realised six-cycle horizon adjudicates. A sharper proportionality of the form $r = \mathcal{O}\!\big(\sigma \cdot \Phi(d/\sqrt{2})\big)$ is a Gaussian-model heuristic that the proof sketch below records as design rationale; its receipt is a measured geometric-rate estimate registered as deferred future work in \Cref{ch:conclusion} under the at-least-fifteen-committed-cycle horizon precondition.
\end{theorem}

\begin{proof}[Proof sketch]
The cycle-boundary update behaves as an asymptotically geometric contraction in expectation under the stationarity assumption stated above. Each cycle admits $n_{\text{new}}$ new entries at the calibration-fold empirical precision $\pi_{\text{store}}$ inherited from \Cref{thm:purity} and overwrites a subset of stored entries' composites under the locked verifier ensemble scoring the updated generator's outputs through the cycle-boundary-refit composite (\Cref{thm:monotone}). The expected reduction in $G_{t}$ per cycle is proportional to the expected number of correctness-flips per stored episode as the cal-fold class-separation $d$ widens under generator advance, which scales as $\Phi(d/\sqrt{2})$ in the Gaussian model, and to the per-cycle admission count, which scales as $\sigma$. Composing these gives an expected per-cycle contraction factor $1-r$ with $r$ as in \Cref{eq:thm-convergence}. The rigorous geometric-envelope claim requires the stochastic-approximation supermartingale machinery of \cite{robbins1951stochastic} and the rate analysis of \cite{polyak1990new} adapted to a bounded gap space; the rigorous supermartingale construction adapted from these works, at the level of a complete proof, is registered as a deferred future-work item in \Cref{ch:conclusion} alongside the at-least-fifteen-committed-cycle horizon precondition that would also unlock the geometric-rate test, and the body claim is therefore stated as an expected-value envelope on the committed-cycle subsequence rather than a sample-path bound on the calendar timeline. The stochastic-envelope refinement of \Cref{cor:stochastic-equilibrium} replaces the calendar-timeline envelope with a Bernoulli-weighted contraction.
\end{proof}

\paragraph{Empirical receipt (bounded-band stationarity on the committed subsequence):} The realised six-cycle trajectory of the calibration-fold storage precision is reported in \Cref{ch:evaluation}: $\pi_{t}$ traces a non-monotone but tightly bounded band of width below $0.03$ across the cycles, with a small downward drift consistent with the per-cycle composite refit and the cycle-four retention abort recorded in \Cref{sec:retention}. Two structural features of the realised horizon limit the strength of the geometric-rate test relative to the stationary expected-value claim of \Cref{eq:thm-convergence}. The first feature is binomial sampling noise: each per-cycle precision reading is a binomial proportion on the cycle's stream-chunk sample, and the resulting per-cycle sampling-noise half-width is comparable in magnitude to the realised cycle-to-cycle gap, so the predicted contraction signal sits at the edge of measurability at the empirical limit of this thesis. The second feature is the cycle-four commit indicator being zero, so the realised committed-cycle subsequence is shorter than the calendar horizon and the per-cycle stationarity of $\sigma$ and $d$ is mildly perturbed by the Phase-aligned composite refit. Under the stochastic-envelope refinement of \Cref{eq:thm-stochastic-envelope}, the realised trajectory satisfies the bounded-distance prediction at the cycle horizon registered for this thesis: the cycle-five reading is the empirical limit against which the theorem is evaluated, the band-width prediction holds, and the implementation regularisers registered at the start of the section account for the gap between this empirical limit and the theoretical asymptote of one. Falsification at the present horizon would require either a monotone-divergent trajectory that exits the band or a monotone-convergent trajectory whose realised gap profile is incompatible with the predicted shape across a horizon long enough to estimate the rate; neither is observed in the six-cycle window. A longer trajectory under fixed thresholds and a fixed composite would be required to test the asymptotic geometric rate directly; this is registered as a future-work item with a fifteen-committed-cycle precondition in \Cref{ch:conclusion}.

\subsection{Asymptotic hallucination floor}
\label{sec:theorem-asymptotic}

The memory-direct tier inherits the storage pool's precision; its asymptotic hallucination rate is therefore upper-bounded by the complement of the pool-precision floor. The full-pipeline decomposition that gathers the memory residual with the architectural false-negative rate is recorded in the scaffolding subsection below because the asymptotic vanishing of the memory residual is not directly observable on a six-cycle horizon.

\begin{corollary}[Tier-1 hallucination floor]
\label{cor:tier1-floor}
Let $\mathrm{CHM}_{1}(t)$ denote the composite hallucination metric of the memory-direct tier at cycle $t$. Under \Cref{thm:purity}, \Cref{thm:convergence}, and the consolidation policy that retains the highest-quality representative within each near-duplicate cluster,
\begin{equation}
\limsup_{t\to\infty}\, \mathrm{CHM}_{1}(t) \;\le\; \max\!\big(1-\pi_{\text{store}},\,1-\pi_{\text{retro}}\big).
\label{eq:thm-asymptotic}
\end{equation}
\end{corollary}

\begin{proof}[Proof sketch]
The memory-direct tier serves the stored answer of the nearest neighbour, so its expected error rate equals the stored pool's expected impurity, $1 - \pi_{t}$. By \Cref{thm:purity}, $\pi_{t} \ge \min(\pi_{\text{store}}, \pi_{\text{retro}})$ for all $t$ that satisfy the preconditions. By \Cref{thm:convergence}, the gap $G_{t} = |\pi_{t} - \pi_{\infty}|$ is bounded above in expectation on the committed-cycle subsequence by a geometric envelope. Combining yields
\(\limsup_{t} \mathrm{CHM}_{1}(t) = 1 - \pi_{\infty} \le 1 - \min(\pi_{\text{store}}, \pi_{\text{retro}}) = \max(1-\pi_{\text{store}}, 1-\pi_{\text{retro}})\),
which is \Cref{eq:thm-asymptotic}. The consolidation policy contributes by ensuring that retrieval at scale does not collapse onto a single contaminated cluster: under the retain-highest-quality rule, the worst entry of any near-duplicate cluster is removed, so the memory-direct tier's effective precision is the maximum across each cluster rather than its mean. The limsup statement in \Cref{eq:thm-asymptotic} tolerates bounded oscillation in the realised pool-precision sequence and should not be read as a pointwise monotone claim.
\end{proof}

\paragraph{Empirical receipt (Tier-1 floor, scope-limited):} The corollary bounds the asymptotic memory-direct hallucination rate by the complement of the stored-pool precision. Two scope conditions must be stated honestly. First, on a held-out evaluation panel constructed to be content-hash disjoint from the training stream, the memory-direct tier fires on a vanishing fraction of queries; the realised trajectory records seven memory-direct dispositions across the six-cycle evaluation horizon at fifteen hundred queries per cycle, all returning the correct stored answer. This rarity is by architectural design rather than a refutation, and it makes the asymptotic trajectory $\mathrm{CHM}_{1}(t)$ statistically unobservable on the realised evaluation distribution. The corollary is therefore a capability bound rather than a trajectory prediction. Second, the conservative ceiling $1-\pi_{\text{store}}\approx 0.36$ derived from the cycle-zero calibration fold overshoots the realised stream-chunk pool complement, which lies in the range $[0.246, 0.277]$ across the realised cycles, and the retroactive prune rate observed at cycle five tightens the binding term to approximately $0.037$. The testable surrogate is the targeted routing probe, in which the memory-direct tier dispatches correctly on verbatim queries to stored episodes at a rate consistent with the calibrated retrieval contract, and every realised Tier-one disposition on the held-out panel produced the correct stored answer at the empirical limit of the trajectory. We therefore present the corollary as a sound pool-level capability bound whose empirical receipt is a probe rather than a trajectory, and we surface the disjointness condition explicitly in the threats to validity.

\subsection{Equilibrium structure and self-correction}
\label{sec:theorem-corollaries}

The corollaries below characterise the equilibrium of the asymptotic statements above and the finite-cycle pruning behaviour that closes the memory-side residual. The self-correction corollary supplies the directional prune-rate receipt on the committed-cycle subsequence; the stochastic-equilibrium corollary supplies the clean empirical pass on the realised commit indicator and the asymmetric-rollback memory invariant. The deployment-cost reading that an explicit half-life corollary would otherwise supply is recovered architecturally through the per-tier cost identity of \Cref{sec:econ-deployment}, with a measured geometric-rate receipt at a horizon of at least fifteen committed cycles registered as future work in \Cref{ch:conclusion}.

\begin{corollary}[Self-correction under improving cal-fold discrimination (locked verifier ensemble, advancing generator)]
\label{cor:self-correction}
Any stored episode whose stored composite was produced when the locked verifier ensemble achieved cal-fold balanced accuracy $\alpha_c$ on the cycle-$c$ generator's outputs is re-scored at every subsequent cycle boundary by the same locked verifier ensemble against the updated generator's outputs and the cycle-boundary-refit composite, yielding the cal-fold balanced accuracies $\alpha_{c+1}, \alpha_{c+2}, \dots$. The verifier ensemble is fixed across the trajectory; the per-cycle changes in $\alpha$ are joint observables of the SIL-improved generator viewed through that fixed lens. Under the assumption $\alpha_{c+k} \ge \alpha_c \ge 1/2$ (the locked verifier discriminates increasingly well on the cycle-$c+k$ generator's cleaner cal-fold distribution because the SIL loop has pushed the generator's distribution onto cleaner support), the conditional probability that the stored episode is correct given that it has survived $k$ retroactive passes converges to $1$ as $k \to \infty$, and any hallucinated episode is pruned within finitely many cycles in expectation.
\end{corollary}

\begin{proof}[Proof sketch]
Treat the indicator $S_k \in \{0,1\}$ that a stored episode survives the $k$-th retroactive pass as a Bernoulli observation through the locked verifier ensemble, with success probability that depends on the latent correctness label $Y \in \{0,1\}$ and on the cal-fold balanced accuracy $\alpha_{c+k}$ at the cycle boundary the pass runs at. Under the locked-verifier-ensemble assumption, $\Pr(S_k=1 \mid Y=1) \ge \alpha_{c+k}$ and $\Pr(S_k=1 \mid Y=0) \le 1-\alpha_{c+k}$, and the survival passes across cycles are conditionally independent given $Y$ because each pass scores the entry against the cycle-boundary-refit composite over the same locked signal-extraction layer. The likelihood-ratio
\[
\Lambda_k := \frac{\Pr(S_1, \dots, S_k \mid Y=1)}{\Pr(S_1, \dots, S_k \mid Y=0)} \;\ge\; \prod_{j=1}^{k} \frac{\alpha_{c+j}}{1 - \alpha_{c+j}}
\]
satisfies $\Lambda_k \to \infty$ as $k \to \infty$ whenever $\alpha_{c+j} \ge \alpha_c > 1/2$ on infinitely many cycles, because each factor exceeds $1$ by a margin bounded below. Bayes' rule then gives $\Pr(Y=1 \mid S_1=\dots=S_k=1) = \Lambda_k \Pr(Y=1) / (\Lambda_k \Pr(Y=1) + \Pr(Y=0)) \to 1$.

For the finite-prune-time claim, fix a hallucinated entry with $Y=0$. The per-cycle prune probability is at least $\alpha_{c+k} - (1-\alpha_{c+k}) = 2\alpha_{c+k} - 1$, bounded below by a strictly positive constant $\delta := 2\alpha_c - 1 > 0$ under the non-decreasing-$\alpha$ premise. The survival probability across $k$ cycles is at most $(1-\delta)^k$, so the expected prune time is at most $\sum_{k=0}^{\infty} (1-\delta)^k = 1/\delta < \infty$ by the geometric-series identity (Borel-Cantelli would give an almost-sure version). The independence-across-cycles assumption is shared with the composite-refit caveat raised on \Cref{thm:monotone}: when the cycle-boundary refit dominates the cal-fold discrimination signal, the conditional-independence premise is approximate and the rate constant $\delta$ is conservative rather than tight.
\end{proof}

\paragraph{Empirical receipt (directional, partial):} The pre-registered receipt was a per-entry survival-cycles distribution across the prune lifecycle. The stored-entry schema persists a storage-cycle field and a retroverified flag but does not carry a passes-survived counter, so the pruned-side histogram of survival counts is deferred to a production-deployment study under the runbook of \Cref{ch:conclusion}. The survivor-side distribution remains recoverable from the cycle-by-cycle memory store snapshots and is reported in \Cref{ch:evaluation} alongside the prune-rate trajectory. The directional receipt on the committed-cycle subsequence is the per-cycle prune rate, which declines monotonically across the four committed retroactive passes from thirteen percent at cycle one to nine percent at cycle two to seven percent at cycle three to four percent at cycle five, consistent with the locked verifier ensemble removing low-quality entries at a decreasing rate as the pool's composition shifts toward higher-confidence survivors and as the SIL-improved generator's outputs widen the cal-fold class separation that the fixed verifier ensemble discriminates against. The decline is partially confounded at the late cycles by fresh-entry dilution from the still-growing pool, so the directional trend is the supported reading rather than a rate estimate. The asymptotic limit $k \to \infty$ is not directly testable at this horizon and is deferred together with the convergence-rate half-life to the longer-horizon future-work item.

\begin{corollary}[Stochastic equilibrium under retention rollback]
\label{cor:stochastic-equilibrium}
Let $P$ denote the multi-modal probe panel registered in \Cref{sec:retention} with floor $\rho_{\min}$. Define the cycle-$t$ commit indicator $X_t \in \{0,1\}$ as the random variable taking value $1$ if the post-fine-tune worst-probe ratio $\min_{p \in P} \rho_p(M_t)$ exceeds $\rho_{\min}$ and $0$ otherwise, with $M_t$ the post-fine-tune model. The commit event depends on the realised training pool composition $Q_t$ (a function of the deferred-buffer state and the stream-chunk yield), the optimizer initialization randomness $s_t$, and the per-probe scoring noise $\varepsilon$. Then:

(i) The expected commit fraction $\pi_{\mathrm{commit}} := \mathbb{E}[X_t]$ satisfies $\pi_{\mathrm{commit}} \in (0, 1]$ on any trajectory where the worst-probe distribution intersects the floor, with $\pi_{\mathrm{commit}} = 1$ only in the degenerate case where the retention guard never fires.

(ii) The convergence envelope of \Cref{thm:convergence} measured on the calendar timeline becomes
\begin{equation}
\mathbb{E}[G_t] \;\le\; (1 - \pi_{\mathrm{commit}} \cdot r)^{t} \, G_0,
\label{eq:thm-stochastic-envelope}
\end{equation}
with $r$ as in \Cref{eq:thm-convergence}; the original $(1-r)^c$ envelope still holds on the committed-cycle subsequence indexed by $c = \sum_{s \le t} X_s$.

(iii) The memory store $|M_t|$ and the per-tier hit-rate trajectory grow monotonically across all cycles because the stream-chunk generate-and-store path executes independently of $X_t$, while the parametric backbone advances only on the committed-cycle subsequence; the cycle-wise architecture therefore decouples deployment-time cost amortisation (memory tier growth) from training-time parametric progression (adapter commits).

(iv) When the worst-probe distribution centres within the probe-noise band $\varepsilon$ of the floor, the trajectory enters a stochastic equilibrium in which commits and rollbacks alternate at rate $\pi_{\mathrm{commit}}$ near the parametric ceiling rather than departing it. The architectural property licensing this equilibrium is the asymmetric rollback registered in \Cref{sec:retention}: every rolled-back cycle preserves the memory store and the deferred buffer exactly as they were at the previous committed cycle's close, so the trajectory can resume parametric advance from any subsequent committed cycle without losing accumulated state.

\paragraph{Scope of the sub-claims:} Sub-claims (i) and (iii) are within-horizon receipts that the realised five-cycle trajectory either satisfies or falsifies directly: the commit-fraction reading is a Bernoulli sample mean on a fully observed indicator sequence, and the memory-monotonicity reading is a structural property of the runner's stream-chunk path that holds at every cycle boundary. Sub-claim (ii) is the calendar-timeline envelope inherited from \Cref{thm:convergence} and is reported as a projected receipt under the same horizon qualifier the parent theorem carries. Sub-claim (iv), the equilibrium-location prediction, is the only sub-claim labelled heuristic in the proof sketch; the proof's stationary-distribution argument is asymptotic, and the realised five-cycle horizon supplies a within-band observation but not a statistical test that distinguishes a stationary equilibrium from a longer-trajectory transient. The rigorous test of sub-claim (iv) requires the at-least-fifteen-committed-cycle horizon registered as future work in \Cref{ch:conclusion} alongside the geometric-rate test of \Cref{thm:convergence}.
\end{corollary}

\begin{proof}[Proof sketch]
Each cycle's commit event is governed by the indicator $X_t = \mathbb{1}[\min_{p \in P} \rho_p(M_t) \ge \rho_{\min}]$, where the worst-probe ratio depends on stochastic factors enumerated above; the indicator is therefore a Bernoulli random variable with success probability $\pi_{\mathrm{commit}}$ that is conditional on the trajectory's history of prior commits (through $Q_t$). Conditional on $X_t = 1$, the per-cycle gap contraction follows \Cref{thm:convergence} at the original rate $r$; conditional on $X_t = 0$, the gap is unchanged because the parameter update is reverted. Taking expectations gives $\mathbb{E}[G_{t+1} \mid G_t] = (1 - \pi_{\mathrm{commit}} \cdot r) \cdot G_t$, and iterating yields \Cref{eq:thm-stochastic-envelope}. The memory-tier growth claim follows from the runner's invariant that the stream-chunk step writes new verified entries to the memory store regardless of whether the cycle's adapter is committed. The equilibrium-location claim is heuristic: when the worst-probe ratio centres within the noise band of the floor, the Bernoulli rate stabilises at a stationary value bounded away from both endpoints, and the trajectory exhibits a bounded distance from the floor across cycles rather than monotone divergence or monotone convergence to a fixed interior point.
\end{proof}

\paragraph{Empirical receipt (stochastic equilibrium):} The receipt is reported sub-claim by sub-claim, separating the within-horizon verifications from the prediction that the five-cycle trajectory cannot adjudicate. Sub-claim (i), the open-interval commit fraction, is verified directly: the realised commit-indicator sequence over the five-cycle trajectory is $X_1{=}1, X_2{=}1, X_3{=}1, X_4{=}0, X_5{=}1$ with worst-probe ratios $1.006, 0.949, 0.976, 0.886, 0.958$ under the locked floor $\rho_{\min} = 0.93$, and the realised commit fraction $\widehat{\pi}_{\mathrm{commit}} = 4/5 = 0.80$ sits strictly in the open interval $(0,1)$. Sub-claim (ii), the calendar-timeline envelope, inherits the bounded-band-versus-deferred-rate split from \Cref{thm:convergence} and is reported under the same horizon qualifier as the parent theorem in \Cref{sec:check-convergence}. Sub-claim (iii), the asymmetric-rollback memory invariant, is also verified directly: the cycle-four abort, triggered by the TriviaQA-test probe collapsing to $0.886$ below the floor, was followed by an immediate restore of $\theta_{\mathrm{prev}}$, but the memory store nonetheless grew across the abort boundary and again across cycle five because the stream-chunk generate-and-store path executed independently of the retention guard, so every cycle, whether committed or rolled back, contributed new verified entries to the memory store while the parametric backbone advanced only on the four committed cycles. Sub-claim (iv), the equilibrium-location prediction, is heuristic at the realised horizon and is reported as such: the worst-probe ratios on cycles two through five all sit within $0.07$ of the floor, qualitatively consistent with the predicted near-floor band, but the five-cycle horizon supplies a within-band observation that is not statistically distinguishable from a longer-trajectory transient at the empirical limit. The rigorous test of sub-claim (iv) requires the at-least-fifteen-committed-cycle horizon registered as future work in \Cref{ch:conclusion}. The corollary's two registered falsifiers, a degenerate commit rate or a memory store that stalls on rolled-back cycles, are both absent within the realised horizon.

\subsection{Mathematical scaffolding}
\label{sec:theorem-scaffolding}

The five identities collected in this subsection underpin the proofs above but do not themselves admit a falsifiable empirical contribution at the empirical limit of the realised trajectory. They are recorded as design-rationale remarks rather than as theorems or corollaries so that the reader can locate the formal machinery in one place without confusing it with a receipt that the evaluation chapter is expected to discharge. Each remark preserves its original label so that cross-references resolved by name continue to point at the same identity.

\begin{remark}[Pool-purity Bayes identity]
\label{thm:bayes-purity}
Let $p_c \in (0,1)$ denote the base-model correctness rate at cycle $c$, and let $\mathrm{TPR}_c$ and $\mathrm{FPR}_c$ denote the storage gate's true-positive and false-positive rates on the labelled purity-validation fold at cycle $c$. Let $P_c$ denote the empirical pool purity. Then by Bayes' rule on the storage event,
\begin{equation}
P_c \;=\; \frac{p_c \cdot \mathrm{TPR}_c}{p_c \cdot \mathrm{TPR}_c \;+\; (1-p_c) \cdot \mathrm{FPR}_c},
\label{eq:bayes-purity}
\end{equation}
and under the symmetric operating-point reduction $\mathrm{TPR}_c = \mathrm{TNR}_c = \alpha_c$ the identity simplifies to $P_c = p_c \alpha_c / (p_c \alpha_c + (1-p_c)(1-\alpha_c))$ with $P_c > p_c$ whenever $\alpha_c > 1/2$. The identity is the mechanism that explains the slack in the calibration-consistent floor of \Cref{thm:purity}: the realised pool purity sits at the Bayes posterior implied by the gate's operating point rather than at the worst-case calibration-fold precision. The labelled purity-validation protocol that would supply $\mathrm{TPR}_c$, $\mathrm{FPR}_c$, and $p_c$ on a held-out labelled fold is registered as a post-trajectory diagnostic and was not consumed by the realised receipt set; the per-cycle residual $|P_c - P_c^{\text{theory}}|$ is therefore listed in \Cref{ch:conclusion} as a future-work diagnostic rather than as a within-horizon empirical check.
\end{remark}

\begin{remark}[Monotonicity of pool purity in base accuracy]
\label{thm:monotone-purification}
Differentiating the symmetric form of the Bayes identity in \Cref{thm:bayes-purity} with respect to $p_c$ under fixed $\alpha > 1/2$ gives
\begin{equation}
\frac{\partial P_c}{\partial p_c} \;=\; \frac{\alpha(1-\alpha)}{\big(p_c \alpha + (1-p_c)(1-\alpha)\big)^{2}} \;>\; 0,
\label{eq:monotone-purification}
\end{equation}
so pool purity is strictly increasing in base accuracy on $(0,1)$ whenever the verifier's balanced accuracy exceeds one half. The identity formalises the intuition that improving the base generator strictly improves the pool the gate produces, with the gain proportional to the gate's discrimination through the $\alpha(1-\alpha)$ factor. The cross-cycle rank-correlation receipt originally registered against this identity is statistically uninformative at the empirical limit of the realised trajectory because the realised range of $p_c$ produces a theoretical $P_c$ span below the per-cycle sampling noise of $P_c$, and because the cycle-boundary recalibration policy lets $\alpha_c$ drift between cycles in violation of the fixed-$\alpha$ premise; the receipt is therefore retired in favour of the calibration-consistent floor reported under \Cref{thm:purity}.
\end{remark}

\begin{remark}[Fixed-point convergence of the Bayes-purification recurrence]
\label{thm:bayes-convergence}
Define the recurrence $P_{c+1} = f(P_c)$ where $f(P) = P \alpha / (P \alpha + (1-P)(1-\alpha))$ under fixed verifier balanced accuracy $\alpha > 1/2$. The map $f: [0,1] \to [0,1]$ is strictly increasing on $(0,1)$ with fixed points only at $\{0, 1\}$. For any $P_0 \in (0, 1)$ the sequence is strictly increasing and bounded above by one, so by the monotone-convergence theorem $P_c \to 1$ as $c \to \infty$. The recurrence is a statement about the pool-purity quantity $P_c$ of \Cref{thm:bayes-purity}, not about the base generator's held-out correctness rate; the fine-tuning step that closes each cycle is a parameter-efficient adapter update with the $L_2$ anchor of \Cref{eq:sil-anchor}, the multi-modal retention guard, and a strict training threshold, none of which implements an idealised posterior update on the held-out generator. The upper fixed point $P^{*} = 1$ is asymptotic and is not observable at the six-cycle empirical limit of this thesis; the implementation regularisers above pin the effective pool-purity ceiling strictly below one, and the realised pool-purity trajectory is therefore expected to fluctuate above the calibration-consistent floor of \Cref{thm:purity} rather than to display strict monotone ascent.
\end{remark}

\begin{remark}[Full-pipeline hallucination decomposition]
\label{thm:asymptotic-elim}
Let $\tau_{1}(t)$ denote the realised Tier-one hit rate on the cycle-$t$ evaluation pool, let $\varepsilon_{\text{arch}}(t)$ denote the joint architectural false-negative rate of the four cascaded storage gates measured at cycle $t$, and let $H_{\text{mem}}(t) = 1 - \pi_{t}$ denote the memory-side hallucination contribution at cycle $t$, where $\pi_{t}$ is the pool purity. By the law of total probability over the tier-routing event, the user-facing hallucination rate satisfies
\begin{equation}
H_t \;\le\; \big(1 - \tau_{1}(t)\big) \cdot \varepsilon_{\text{arch}}(t) \;+\; \tau_{1}(t) \cdot H_{\text{mem}}(t).
\label{eq:asymptotic-elim}
\end{equation}
The pool-purity term is bounded below by the calibration-consistent floor of \Cref{thm:purity}, and the equilibrium residual at the implementation's effective limit decomposes as $H_{\infty} \approx (1 - \tau_{1}^{*}(D)) \cdot (1 - x(\lvert\theta\rvert, C))$ where $x(\lvert\theta\rvert, C)$ is the asymptotic correctness rate of the base generator with parameter count $\lvert\theta\rvert$ over retrieval corpus $C$. The vanishing-memory-residual asymptote $H_{\text{mem}}(t) \to 0$ is an infinite-horizon claim that is unobservable at the empirical limit of the trajectory for two structural reasons: the evaluation pool is content-hash disjoint from the training stream by design, pinning $\tau_{1}(t)$ at the noise floor on the realised horizon; and the retention-guard rollback, the parametric capacity of the base generator, and the regularisation anchor pin the pool-purity ceiling strictly below one. The deployment-mode setting registered in \Cref{ch:conclusion} under the production runbook, in which repeat queries from the served distribution re-enter the routing layer and the cycle horizon extends to the calendar trigger, is the regime in which the asymptotic statements of this decomposition become testable. The tier-partition decomposition is routing-agnostic in the partition step; the $\varepsilon_{\text{arch}}(t)$ term is measured against the realised post-classifier tier distribution of \Cref{sec:ruc}, so the corpus-conditional residual $1 - x(\lvert\theta\rvert, C)$ at the asymptote is the corpus-conditional bound on the retrieval-augmented branch of the post-classifier split, with the zero-shot branch's residual bounded by the parametric ceiling of the fine-tuned generator rather than by the corpus floor.
\end{remark}

\begin{remark}[Corpus-bounded hallucination floor]
\label{cor:corpus-floor}
Let $K(C) \subseteq D$ denote the set of queries whose true answer is supportable by retrieved evidence from the corpus $C$. Under the grounded-reasoning convention, no system that asserts only claims supportable by $C$ can achieve a hallucination rate strictly below $\Pr_{D}[q \notin K(C)]$, so the equilibrium residual of \Cref{thm:asymptotic-elim} is bounded below by an information-theoretic floor determined by the corpus rather than by the architecture. The full set $K(C)$ is not directly observable, and the pre-registered proxy diagnostic, which would identify evaluation-fold queries whose top-$k$ retrieval does not contain the gold-answer span and report the composite hallucination metric on that subset, was not consumed by the realised receipt set because the rectangular-record evaluation harness intentionally filters list-typed verifier fields and the proxy reconstruction therefore requires re-running retrieval against the dense passage index for every evaluation question across all six cycles. The proxy reconstruction is registered as a deferred receipt in \Cref{ch:conclusion}; the corpus-bounded floor itself enters the present chapter as an interpretive limit on the architectural equilibrium rather than as a measured quantity at the empirical limit of the trajectory.
\end{remark}


\section{Chapter signpost}
\label{sec:ch4-signpost}

The methodology specified in this chapter is intentionally complete: every component, every threshold, and every fit-time procedure is registered before the empirical record is consulted. \Cref{ch:evaluation} reads this specification verbatim against the realised cycle trajectory, the external baseline panel, and the retrieval-utility-classifier on/off architectural ablation, and reports each registered hypothesis as supported, partially supported, or unsupported with the deciding statistic. The load-bearing theorems and corollaries of \Cref{sec:theory} are checked there one by one against their empirical receipts at the empirical limit of the realised trajectory, so the reader can locate each architectural claim in the specification chapter and find the corresponding receipt in the evaluation chapter without ambiguity.
